%% Question bank chapter file - generated by bank_tools.py build (do not edit while the bank is still being built:
%% it is regenerated from the per-test data in Question Bank\data). Every question is one bankq block:
%%   \begin{bankq}[id=..., source={...}, paper=.., qno=.., type=..., subtopic={...}, answer={...}, answer_src=key|solved,
%%                 working={...}, difficulty=1-10, diff_parts={K S I C F}, diff_note={...},
%%                 group=<passage id or empty>, also_in={duplicate copies in other tests}]
%% difficulty = min(10, K+S+I+C+F-1): K concepts 1-3, S steps 1-4, I insight 0-2, C computation 0-2, F format 0-1
%%   ...question LaTeX, options, figures (\bfig / \bfigi -> img/)...
%%   \end{bankq}
%% Passages shared by several questions are bankpassage blocks; their questions carry group=<passage id>.
\documentclass[10pt]{article}
\usepackage[a4paper,margin=1.8cm]{geometry}
\usepackage[T1]{fontenc}
\usepackage{amsmath,amssymb,mathtools}
\usepackage{graphicx}
\usepackage[export]{adjustbox}
\usepackage{xcolor,needspace,keyval}
\usepackage{enumitem}
\usepackage{booktabs,array,tabularx}
\usepackage[hidelinks]{hyperref}
\setlength{\parindent}{0pt}
\setlength{\parskip}{3pt}
\makeatletter
\define@key{bq}{id}{\def\bq@id{#1}}\define@key{bq}{source}{\def\bq@src{#1}}\define@key{bq}{paper}{\def\bq@paper{#1}}
\define@key{bq}{qno}{\def\bq@qno{#1}}\define@key{bq}{type}{\def\bq@type{#1}}\define@key{bq}{subtopic}{\def\bq@sub{#1}}
\define@key{bq}{answer}{\def\bq@ans{#1}}\define@key{bq}{answer_src}{\def\bq@asrc{#1}}\define@key{bq}{working}{\def\bq@work{#1}}
\define@key{bq}{group}{\def\bq@group{#1}}\define@key{bq}{also_in}{\def\bq@also{#1}}
\define@key{bq}{difficulty}{\def\bq@diff{#1}}\define@key{bq}{diff_parts}{\def\bq@dparts{#1}}\define@key{bq}{diff_note}{\def\bq@dnote{#1}}
\newenvironment{bankq}[1][]{\def\bq@work{}\def\bq@also{}\def\bq@group{}\def\bq@diff{?}\def\bq@dparts{}\def\bq@dnote{}\setkeys{bq}{#1}%
  \par\needspace{4\baselineskip}\medskip\noindent{\small\textbf{\bq@id}\hfill\color{gray}[\bq@src, Paper \bq@paper, Q\bq@qno]
  \bq@type\ $\cdot$ \bq@sub\ $\cdot$ \textbf{D\bq@diff}}\par\nopagebreak\relax}
  {\par{\small\color{teal}Answer: \textbf{\bq@ans}\ifx\bq@work\@empty\else\ (solved: \bq@work)\fi
  \ifx\bq@also\@empty\else\par\color{gray}Also in: \bq@also\fi}\par}
\define@key{bp}{id}{\def\bp@id{#1}}\define@key{bp}{source}{\def\bp@src{#1}}\define@key{bp}{questions}{\def\bp@qs{#1}}
\newenvironment{bankpassage}[1][]{\setkeys{bp}{#1}\par\needspace{4\baselineskip}\bigskip\noindent
  {\small\textbf{Passage \bp@id}\hfill\color{gray}[\bp@src; questions \bp@qs]}\par\nopagebreak\relax\itshape}{\par}
\makeatother
\newcommand{\bfig}[2]{\par\smallskip{\centering\includegraphics[width=#1pt,max width=0.62\linewidth,max height=0.32\textheight]{img/#2}\par}\smallskip}
\newcommand{\bfigi}[2]{\includegraphics[width=#1pt,max width=\linewidth,max height=0.32\textheight,valign=c]{img/#2}}
\newcommand{\bmissing}[1]{\textit{[figure missing in the source conversion: \detokenize{#1}]}}
\newenvironment{qoptions}{\begin{itemize}[leftmargin=1.8em,itemsep=1pt,topsep=2pt,label={}]}{\end{itemize}}
\newcommand{\opt}[2]{\item[(#1)] #2}
\providecommand{\fwstart}{}\providecommand{\fwend}{}\providecommand{\qfullsplit}{}
\newcolumntype{L}[1]{>{\raggedright\arraybackslash}p{#1}}
\newcolumntype{C}[1]{>{\centering\arraybackslash}p{#1}}
\newcolumntype{R}[1]{>{\raggedleft\arraybackslash}p{#1}}
\providecommand{\sectionhdr}[1]{}\providecommand{\subsectionhdr}[1]{}


\begin{document}
{\Large\bfseries GOC --- Question Bank}\par\smallskip
{\small Chemistry chapter C14. Questions from all converted test papers, grouped by source series. Answers: printed key where it matches the paper, otherwise solved (answer\_src=solved, with working). Exact repeats across papers appear once, with also\_in listing the other copies.}\par\medskip\hrule

\section*{Score-I papers}
\begin{bankq}[id=SI06-P1-CHEM-S1-Q05, source={Score-I - Test 06 (28 Jan 2022)}, paper=1, qno=5, type=single, subtopic={Reactive intermediates}, answer={C}, answer_src=key, difficulty=2, diff_parts={K1 S1 I1 C0 F0}, diff_note={concepts: stability order of drawn species | steps: resonance, hyperconjugation and inductive effects; rank a b c}]
The correct stability order of the following species is.
\bfig{283.1}{SI06_p1_chem_s1_q05_1.png}

\begin{qoptions}
\opt{A}{c $<$ a $<$ b}
\opt{B}{c = b $<$ a}
\opt{C}{c $<$ a = b}
\opt{D}{a = b = c}
\end{qoptions}
\end{bankq}
\begin{bankq}[id=SI06-P1-CHEM-S1-Q06, source={Score-I - Test 06 (28 Jan 2022)}, paper=1, qno=6, type=single, subtopic={Electronic effects}, answer={B}, answer_src=key, difficulty=3, diff_parts={K1 S2 I1 C0 F0}, diff_note={concepts: incorrect statement about enol content | steps: enolisation favoured by conjugation and hydrogen bonding; compare compounds; find incorrect}]
Which of the following is incorrect for Enol \% ?
\begin{qoptions}
\opt{A}{\bfigi{124.2}{SI06_p1_chem_s1_q06_A.png}}
\opt{B}{\bfigi{95.4}{SI06_p1_chem_s1_q06_B.png}}
\opt{C}{\bfigi{98.9}{SI06_p1_chem_s1_q06_C.png}}
\opt{D}{\bfigi{87.5}{SI06_p1_chem_s1_q06_D.png}}
\end{qoptions}
\end{bankq}
\begin{bankq}[id=SI07-P1-CHEM-S2-Q05, source={Score-I - Test 07 (04 Feb 2022)}, paper=1, qno=5, type=numerical, subtopic={Acidity and basicity}, answer={4}, answer_src=key, difficulty=5, diff_parts={K1 S2 I1 C1 F1}, diff_note={concepts: count compounds more acidic than carbonic acid and less acidic than a reference | steps: pKa values of phenol, carboxylic acids, water; compare with 6.4; count}]
Compounds having pKa value less than that of carbonic acid (i.e., H$_2$CO$_3$ 6.4 approx.) is X and greater than that of carbonic acid is Y. Then find out $\dfrac{\mathrm X-\mathrm Y}{3}$ =?

(i) \bfigi{47.7}{SI07_p1_chem_s2_q05_i.png}

(ii) \bfigi{53.7}{SI07_p1_chem_s2_q05_ii.png}

(iii) \bfigi{73.7}{SI07_p1_chem_s2_q05_iii.png}

(iv) \bfigi{75.2}{SI07_p1_chem_s2_q05_iv.png}

(v) \bfigi{57.8}{SI07_p1_chem_s2_q05_v.png}

(vi) \bfigi{54.6}{SI07_p1_chem_s2_q05_vi.png}

(vii) \bfigi{54.8}{SI07_p1_chem_s2_q05_vii.png}

(viii) \bfigi{25.0}{SI07_p1_chem_s2_q05_viii.png}

(ix) \bfigi{56.8}{SI07_p1_chem_s2_q05_ix.png}

(x) \bfigi{53.7}{SI07_p1_chem_s2_q05_x.png}
\end{bankq}
\begin{bankq}[id=SI07-P1-CHEM-S3-Q01, source={Score-I - Test 07 (04 Feb 2022)}, paper=1, qno=1, type=numerical, subtopic={Electronic effects}, answer={3.14}, answer_src=key, difficulty=6, diff_parts={K2 S2 I1 C1 F1}, diff_note={concepts: count groups showing resonance, a phenyl ring | steps: groups with lone pairs or pi bonds directly attached; allyl no; vinyl yes; count}]
How many groups will show resonance with phenyl ring when attached with it?

(a) $-\mathrm{CH_2-CH=CH_2}$ \quad (b) $-\mathrm{CH_2-O-\overset{\mathrm O}{\overset{\|}{C}}-NH_2}$

(c) $-\overset{\oplus}{\mathrm N}\mathrm H_3$ \quad (d) $-\mathrm{N=O}$

(e) $-\mathrm{Cl}$ \quad (f) $-\overset{\oplus}{\mathrm C}\mathrm H_2$

(g) \bfigi{36.3}{SI07_p1_chem_s3_q01_g.png} \quad (h) $-\mathrm{BH_2}$
\end{bankq}
\begin{bankq}[id=SI07-P1-CHEM-S3-Q02, source={Score-I - Test 07 (04 Feb 2022)}, paper=1, qno=2, type=numerical, subtopic={Electronic effects}, answer={66.66 to 66.67}, answer_src=key, difficulty=6, diff_parts={K2 S2 I1 C1 F1}, diff_note={concepts: count compounds, C=C shorter than propene | steps: hyperconjugation and hybridisation effects on bond length; compare each; count}]
How many compounds given below have C = C bond length shorter than $\mathrm{H_2C\underset{\uparrow}{=}CH-CH_3}$.

(a) \bfigi{32.5}{SI07_p1_chem_s3_q02_a.png} \quad (b) \bfigi{26.3}{SI07_p1_chem_s3_q02_b.png}

(c) $\mathrm{O_2N-CH\underset{\uparrow}{=}CH-OEt}$ \quad (d) \bfigi{26.5}{SI07_p1_chem_s3_q02_d.png}

(e) $\mathrm{CH_3-CH_2-CH\underset{\uparrow}{=}CH_2}$

(f) $\mathrm{H_3C-CH\underset{\uparrow}{=}CH-CH_3}$

(g) $\mathrm{(CH_3)_2C\underset{\uparrow}{=}C(CH_3)_2}$
\end{bankq}
\begin{bankq}[id=SI08-P2-CHEM-S1-Q05, source={Score-I - Test 08 (09 Feb 2022)}, paper=2, qno=5, type=multi, subtopic={Reactive intermediates}, answer={A,B,D}, answer_src=key, difficulty=4, diff_parts={K1 S2 I1 C0 F1}, diff_note={concepts: stability order of drawn species | steps: carbocation stability or resonance; compare; check orders}]
Which of the given following stability order is correct.
\begin{qoptions}
\opt{A}{\bfigi{150.1}{SI08_p2_chem_s1_q05_A.png}}
\opt{B}{\bfigi{115.3}{SI08_p2_chem_s1_q05_B.png}}
\opt{C}{\bfigi{162.1}{SI08_p2_chem_s1_q05_C.png}}
\opt{D}{\bfigi{160.9}{SI08_p2_chem_s1_q05_D.png}}
\end{qoptions}
\end{bankq}
\begin{bankq}[id=SI10-P1-CHEM-S3-Q01, source={Score-I - Test 10 (24 Feb 2022)}, paper=1, qno=1, type=matching, subtopic={Reactive intermediates}, answer={A$\to$P,R; B$\to$R,S; C$\to$R,S; D$\to$Q,R,T}, answer_src=solved, working={A: tertiary alcohol dehydration, E1 through a classical carbocation. B: acid-catalysed hydration of the propargylic alcohol, vinyl carbocation then rearrangement (Meyer-Schuster). C: dienone-phenol rearrangement, carbocation and methyl shift. D: ICN adds across the ring, iodonium (non-classical) ion, trans addition gives a racemic pair.}, difficulty=6, diff_parts={K2 S3 I1 C0 F1}, diff_note={concepts: E1 and carbocation chemistry, stereochemistry of addition | steps: A dehydration of the tertiary alcohol is E1 with a classical carbocation; B hydration of the propargyl alcohol goes through a vinyl cation and rearranges; C dienone-phenol rearrangement has a carbocation and a methyl shift; D iodonium gives trans addition and a racemic pair; classical against non-classical cation; assign the five labels | I1: recognising the two named rearrangements}]
Match the entries of Column-I with Column-II.
\begin{center}\footnotesize
\begin{tabular}{|p{0.6cm}|p{6.6cm}|}
\hline
\multicolumn{2}{|c|}{Column - I (Chemical reaction)}\\ \hline
(A) & \bfigi{127.6}{SI10_p1_chem_s3_pre_A.png}\\ \hline
(B) & \bfigi{156.2}{SI10_p1_chem_s3_pre_B.png}\\ \hline
(C) & \bfigi{102.3}{SI10_p1_chem_s3_pre_C.png}\\ \hline
(D) & \bfigi{114.1}{SI10_p1_chem_s3_pre_D.png}\\ \hline
\end{tabular}
\end{center}
\begin{center}\footnotesize
\begin{tabular}{|p{0.6cm}|p{6.6cm}|}
\hline
\multicolumn{2}{|c|}{Column - II (Type of reaction/intermediate involve)}\\ \hline
(P) & Elimination reaction\\ \hline
(Q) & Addition reaction\\ \hline
(R) & Carbocation intermediate is formed [Classical /Non classical]\\ \hline
(S) & Rearrangment reaction\\ \hline
(T) & Racemic mixture is formed\\ \hline
\end{tabular}
\end{center}
\end{bankq}
\begin{bankq}[id=SI11-P1-CHEM-S2-Q08, source={Score-I - Test 11 (27 Feb 2022)}, paper=1, qno=8, type=multi, subtopic={Acidity and basicity}, answer={A,B,C}, answer_src=key, difficulty=7, diff_parts={K2 S3 I1 C1 F1}, diff_note={concepts: acidity, conjugate base properties of compounds I to V | steps: delocalisation stabilises conjugate bases; aromatic conjugate base of cyclopentadiene type; nitro substituent increases acidity; order of acidity; check each statement}]
With respect to the compounds \textbf{I}--\textbf{V}, choose the correct statement(s).
\bfig{245.2}{SI11_p1_chem_s2_q08_1.png}
\begin{qoptions}
\opt{A}{The acidity of compound \textbf{I} is due to delocalization in the conjugate base.}
\opt{B}{The conjugate base of compound \textbf{IV} is aromatic.}
\opt{C}{Compound \textbf{II} becomes more acidic, when it has a $-\mathrm{NO_2}$ substituent.}
\opt{D}{The acidity of compounds follows the order \textbf{I} $>$ \textbf{IV} $>$ \textbf{V} $>$ \textbf{II} $>$ \textbf{III}.}
\end{qoptions}
\end{bankq}
\begin{bankq}[id=SI11-P2-CHEM-S1-Q05, source={Score-I - Test 11 (27 Feb 2022)}, paper=2, qno=5, type=multi, subtopic={Acidity and basicity}, answer={C,D}, answer_src=key, difficulty=6, diff_parts={K2 S2 I1 C1 F1}, diff_note={concepts: basicity order, resonance of four nitrogen compounds | steps: compare lone pair availability and resonance in each; order of basicity; pKb differences; check statements}]
Consider the following four compounds \textbf I, \textbf{II}, \textbf{III}, and \textbf{IV}.
\bfig{245.2}{SI11_p2_chem_s1_q05_1.png}
Choose the correct statement(s).
\begin{qoptions}
\opt{A}{The order of basicity is \textbf{II} $>$ \textbf I $>$ \textbf{III} $>$ \textbf{IV}.}
\opt{B}{The magnitude of $\mathrm{p}K_b$ difference between \textbf I and \textbf{II} is more than that between \textbf{III} and \textbf{IV}.}
\opt{C}{Resonance effect is more in \textbf{III} than in \textbf{IV}.}
\opt{D}{Steric effect makes compound \textbf{IV} more basic than \textbf{III}.}
\end{qoptions}
\end{bankq}
\begin{bankpassage}[id=SI12-P2-CHEM-S1-G1-2, source={Score-I - Test 12 (09 Mar 2022)}, questions={1-2}]
\begin{center}\small
\begin{tabular}{|p{0.8cm}|p{6.2cm}|p{0.8cm}|p{4.0cm}|p{0.8cm}|p{3.0cm}|}
\hline
 & \centering Column-I & & \centering Column-II & & \centering Column-III\tabularnewline
\hline
(I) & \bfigi{65.6}{SI12_p2_chem_s1_para1_2_1.png}\rule{0pt}{20pt} & (i) & Carbene intermediate & (a) & Haffmann product\\
\hline
(II) & \bfigi{107.0}{SI12_p2_chem_s1_para1_2_2.png}\rule{0pt}{40pt} & (ii) & Carbocation intermediate & (b) & Aromatic compound\\
\hline
(III) & \bfigi{107.0}{SI12_p2_chem_s1_para1_2_3.png}\rule{0pt}{32pt} & (iii) & Transition state & (c) & 2-methyl But-2-ene\\
\hline
(IV) & \bfigi{57.3}{SI12_p2_chem_s1_para1_2_4.png}\rule{0pt}{36pt} & (iv) & Free radical intermediate & (d) & cis alkene\\
\hline
\end{tabular}
\end{center}
\end{bankpassage}
\begin{bankq}[id=SI12-P2-CHEM-S1-Q01, source={Score-I - Test 12 (09 Mar 2022)}, paper=2, qno=1, type=single, subtopic={Reactive intermediates}, answer={B}, answer_src=key, difficulty=2, diff_parts={K2 S1 I0 C0 F0}, diff_note={concepts: carbocation intermediates, rearrangement | steps: identify the intermediate from the reaction; match the types and the outcomes}, group=SI12-P2-CHEM-S1-G1-2]
Correct set for following reaction
\begin{qoptions}
\opt{A}{I, ii, d}
\opt{B}{I, iii, a}
\opt{C}{I, i, c}
\opt{D}{I, iii, b}
\end{qoptions}
\end{bankq}
\begin{bankq}[id=SI12-P2-CHEM-S1-Q02, source={Score-I - Test 12 (09 Mar 2022)}, paper=2, qno=2, type=single, subtopic={Reactive intermediates}, answer={C}, answer_src=key, difficulty=2, diff_parts={K2 S1 I0 C0 F0}, diff_note={concepts: carbocation intermediates, rearrangement | steps: identify the intermediate from the reaction; match the types and the outcomes}, group=SI12-P2-CHEM-S1-G1-2]
Correct set for following reaction :
\begin{qoptions}
\opt{A}{III, ii, c}
\opt{B}{IV, iii, d}
\opt{C}{II, iii, b}
\opt{D}{All}
\end{qoptions}

\noindent
\end{bankq}
\begin{bankq}[id=SI17-P2-CHEM-S3-Q05, source={Score-I - Test 17 (17 Apr 2022)}, paper=2, qno=5, type=numerical, subtopic={Reactive intermediates}, answer={8}, answer_src=key, difficulty=7, diff_parts={K2 S3 I1 C1 F1}, diff_note={concepts: difference of degree of unsaturation, rearrangements in a sequence | steps: follow carbocation rearrangements through steps; product Z structure; degree of unsaturation; number of rearrangements; subtract}]
Observe the following reaction sequence
\bfig{299.5}{SI17_p2_chem_s3_q05_1.png}

Degree of unsaturation of product (Z) = X

Total number of rearrangement in reaction sequence from starting compound to Z = Y

Find X $-$ Y = ?
\end{bankq}
\begin{bankq}[id=SI18-P1-CHEM-S2-Q08, source={Score-I - Test 18 (24 Apr 2022)}, paper=1, qno=8, type=multi, subtopic={Reactive intermediates}, answer={A,B,C}, answer_src=key, difficulty=6, diff_parts={K1 S2 I2 C1 F1}, diff_note={concepts: concepts linked to the drawn reaction mechanism | steps: identify the reaction involving carbocation formation; rearrangement by methyl shift; Friedel Crafts type ring closure; check options | I: track the carbocation through shifts}]
Which of the following are linked to the reaction ?
\bfig{152.1}{SI18_p1_chem_s2_q08_1.png}
\begin{qoptions}
\opt{A}{carbocation rearrangement}
\opt{B}{\bfigi{64.4}{SI18_p1_chem_s2_q08_B.png}}
\opt{C}{Friedel-Crafts reaction}
\opt{D}{Methyl shift}
\end{qoptions}
\end{bankq}
\begin{bankq}[id=SI22-P1-CHEM-S1-Q02, source={Score-I - Test 22 (04 Dec 2022)}, paper=1, qno=2, type=multi, subtopic={Electronic effects}, answer={B,C,D}, answer_src=key, difficulty=5, diff_parts={K1 S2 I1 C1 F1}, diff_note={concepts: species in which the charge is involved in resonance | steps: check whether the charged atom is conjugated with a pi system or lone pair; benzyl and allyl cations and anions involve resonance; sp3 carbanions near sigma bonds do not; select}]
In which of the following species, the positive or negative charge is involved in resonance.
\begin{qoptions}
\opt{A}{\bfigi{25.2}{SI22_p1_chem_s1_q02_A.png}}
\opt{B}{\bfigi{48.7}{SI22_p1_chem_s1_q02_B.png}}
\opt{C}{\bfigi{40.9}{SI22_p1_chem_s1_q02_C.png}}
\opt{D}{\bfigi{26.2}{SI22_p1_chem_s1_q02_D.png}}
\end{qoptions}
\end{bankq}
\begin{bankpassage}[id=SI22-P1-CHEM-S2-G9-10, source={Score-I - Test 22 (04 Dec 2022)}, questions={9-10}]
In the following, List-I contains some organic compounds, and List-II contains properties.
\par\smallskip
\begin{center}\small
\begin{tabular}{|c|c|c|p{4.2cm}|}
\hline
 & List-I & & List-II\\
\hline
(I) & \bfigi{59.0}{SI22_p1_chem_s2_para9_10_1.png}\rule{0pt}{26pt} & (P) & Aromatic compound\\
\hline
(II) & \bfigi{54.1}{SI22_p1_chem_s2_para9_10_2.png}\rule{0pt}{26pt} & (Q) & More acidic than C$_2$H$_5$OH\\
\hline
(III) & \bfigi{44.1}{SI22_p1_chem_s2_para9_10_3.png}\rule{0pt}{26pt} & (R) & Undergoes addition reaction with electrophiles\\
\hline
(IV) & \bfigi{58.3}{SI22_p1_chem_s2_para9_10_4.png}\rule{0pt}{26pt} & (S) & Undergoes substitution reaction with odd electron species\\
\hline
 & & (T) & Produces gas with Grignard reagent.\\
\hline
 & & (U) & Soluble in NaOH\\
\hline
\end{tabular}
\end{center}
\end{bankpassage}
\begin{bankq}[id=SI22-P1-CHEM-S2-Q09, source={Score-I - Test 22 (04 Dec 2022)}, paper=1, qno=9, type=single, subtopic={Acidity and basicity}, answer={B}, answer_src=key, difficulty=6, diff_parts={K1 S3 I2 C1 F0}, diff_note={concepts: properties of three organic compounds from list I matched with tests in list II | steps: identify each drawn compound; apply tests like acidity with bicarbonate; Tollens reagent; iodoform; Lucas; decide which properties each shows; choose the correct set for compounds II and III | I: several properties tested for each compound}, group=SI22-P1-CHEM-S2-G9-10]
Which of the following options has the correct combination considering List-I and List-II.
\begin{qoptions}
\opt{A}{(II) -- P, Q, R, S, T, U}
\opt{B}{(III) -- P, R}
\opt{C}{(II) -- P, Q, R, S, T}
\opt{D}{(III) -- R, S, T}
\end{qoptions}
\end{bankq}
\begin{bankq}[id=SI22-P1-CHEM-S2-Q10, source={Score-I - Test 22 (04 Dec 2022)}, paper=1, qno=10, type=single, subtopic={Acidity and basicity}, answer={D}, answer_src=key, difficulty=5, diff_parts={K1 S2 I2 C1 F0}, diff_note={concepts: properties of compounds I and IV matched with tests | steps: same list of compounds; identify compound I and IV functional groups; match with properties P S T U; the correct set | I: need to reason through four properties for the two compounds}, group=SI22-P1-CHEM-S2-G9-10]
Which of the following options has the correct combination considering List-I and List-II.
\begin{qoptions}
\opt{A}{(I) -- P, S, U}
\opt{B}{(IV) -- P, T}
\opt{C}{(I) -- P, S, T}
\opt{D}{(IV) -- S}
\end{qoptions}
\end{bankq}
\begin{bankq}[id=SI22-P1-CHEM-S3-Q02, source={Score-I - Test 22 (04 Dec 2022)}, paper=1, qno=2, type=numerical, subtopic={Reactive intermediates}, answer={7.00}, answer_src=key, difficulty=5, diff_parts={K1 S2 I1 C1 F1}, diff_note={concepts: count electrophilic species | steps: electron deficient cations and Lewis acids are electrophiles; SO3 acylium diazonium chlorine; permanganate oxidant; count}]
How many of the following species are electrophilic in nature.\\
\bfigi{85.4}{SI22_p1_chem_s3_q02_1.png}, SO$_3$, CH$_3$--C$\equiv$O$^+$, Ph--$\overset{+}{\text N}\equiv$N, Cl$_2$, MnO$_4^-$
\end{bankq}
\begin{bankq}[id=SI31-P1-CHEM-S1-Q07, source={Score-I - Test 31 (29 Dec 2022)}, paper=1, qno=7, type=multi, subtopic={Electronic effects}, answer={A,B,C}, answer_src=key, difficulty=7, diff_parts={K1 S3 I2 C1 F1}, diff_note={concepts: correct orders for electrophilic addition radical stability basicity acidity | steps: HCl addition rate by carbocation stability; radical stability with allyl and benzyl; basic strength in water of amines; acidic strength of acids; check each | I: four different property orders}]
Correct option among the following is/are.
\begin{qoptions}
\opt{A}{Ease of HCl addition :

$>$C = CH$_2$ $>$ CH$_3$CH = CH$_2$ $>$ HC$\equiv$CH}
\opt{B}{Order of Free Radial Stability\\$\mathrm{CH_2=CH-\dot{C}H_2>(CH_3)_3\dot{C}>CH_3\dot{C}H_2>\dot{P}h}$}
\opt{C}{Order of Basic strength in aq. medium :

Me$_2$NH $>$ MeNH$_2$ $>$ Me$_3$N $>$ NH$_3$}
\opt{D}{Order of Acidic strength: HCOOH $>$ ClCH$_2$COOH $>$ HOOCCH$_2$--COOH $>$ $\mathrm{(CH_3)_3\overset{\oplus}{N}CH_2COOH}$}
\end{qoptions}
\end{bankq}

\section*{Score-II papers}
\begin{bankq}[id=SII05-P1-CHEM-S1-Q10, source={Score-II - Test 05 (22 Mar 2022)}, paper=1, qno=10, type=single, subtopic={Electronic effects}, answer={D}, answer_src=key, difficulty=2, diff_parts={K1 S1 I1 C0 F0}, diff_note={concepts: least stable resonance structure | steps: charge separation, octet incomplete, electronegativity of the atom bearing the charge; compare the structures}]
Among the following the least stable resonance structure is :
\begin{qoptions}
\opt{A}{\bfigi{92.0}{SII05_p1_chem_s1_q10_A.png}}
\opt{B}{\bfigi{84.6}{SII05_p1_chem_s1_q10_B.png}}
\opt{C}{\bfigi{91.0}{SII05_p1_chem_s1_q10_C.png}}
\opt{D}{\bfigi{91.0}{SII05_p1_chem_s1_q10_D.png}}
\end{qoptions}
\end{bankq}
\begin{bankq}[id=SII15-P1-CHEM-S1-Q20, source={Score-II - Test 15 (09 Feb 2024)}, paper=1, qno=20, type=single, subtopic={Acidity and basicity}, answer={C}, answer_src=key, difficulty=3, diff_parts={K1 S2 I1 C0 F0}, diff_note={concepts: order of compounds against indicated property | steps: stability of Pb IV chloride; magnetic moments; compare}]
Which of the following order of compound is correct against to their indicated property :-
\begin{qoptions}
\opt{A}{PbCl$_2$ $<$ PbCl$_4$ (Stability)}
\opt{B}{FeCl$_2$ $>$ FeCl$_3$ (Covalent nature)}
\opt{C}{LiF $>$ NaF (Thermal stability)}
\opt{D}{N$_2$O $>$ NO (Acidic nature)}
\end{qoptions}
\end{bankq}
\begin{bankq}[id=SII15-P1-CHEM-S2-Q10, source={Score-II - Test 15 (09 Feb 2024)}, paper=1, qno=10, type=numerical, subtopic={Electronic effects}, answer={9}, answer_src=key, difficulty=4, diff_parts={K1 S1 I1 C1 F1}, diff_note={concepts: count groups that donate electron density to the benzene ring | steps: resonance donating groups; count}]
How many of the following groups are capable of donating electron density to the benzene ring during the Wheland intermediate, being ortho / para directors :\\
(a) $-$Cl \hfill (b) $-$NO \hfill (c) NO$_2$\\
(d) $-$CCl$_3$ \hfill (e) $-$CH=CH$_2$\\
(f) $-$CH = CH $-$ Cl. \hfill (g) $-$CH = CH $-$ NO$_2$
\bfig{147.0}{SII15_p1_chem_s2_q10_1.png}
\end{bankq}
\begin{bankq}[id=SII16-P1-CHEM-S1-Q02, source={Score-II - Test 16 (09 Feb 2024)}, paper=1, qno=2, type=multi, subtopic={Acidity and basicity}, answer={A,B,D}, answer_src=key, difficulty=6, diff_parts={K1 S3 I1 C1 F1}, diff_note={concepts: four ionic equilibria statements | steps: fraction ionised of a weak base; pH of mixing an acid and a base of equal strength; indicator ionisation percentage from Henderson equation; common-ion solubility of AgCl in MgCl2; verify each of four}]
Identify correct statement(s) from the following given options :
\begin{qoptions}
\opt{A}{The expression for \% ionisation of a dilute solution of a monoacidic weak base in aqueous solution is $\dfrac{100}{1+10^{(pK_b-pOH)}}$}
\opt{B}{If equal volumes of following solutions are mixed, then the pH of resulting solution will be average value of pH of two given solutions.\\
Solution I : pH = 5 ($CH_3COOH$)(aq) \& Solution-II : pH = 9 ($NH_3$(aq)),\\
Given : $[K_a(CH_3COOH)=K_b(NH_3(aq))]$}
\opt{C}{An acidic indicator HIn (colourless) will be visible only when its ionised form (pink in colour) is 25 \% by mole or more of total indicator in a solution. Suppose HIn ($pK_a=9.0$) is added to a solution having pH = 9.6 then mole \% of ionised form of HIn in solution is 90\%, (Take log 2 = 0.3)}
\opt{D}{Molar solubility of AgCl ($K_{SP}=10^{-12}$) in 0.2 M $MgCl_2$ aqueous solution is $2.5\times10^{-12}$ M.}
\end{qoptions}
\end{bankq}
\begin{bankq}[id=SII18-P1-CHEM-S1-Q06, source={Score-II - Test 18 (13 Feb 2024)}, paper=1, qno=6, type=single, subtopic={Acidity and basicity}, answer={D}, answer_src=key, difficulty=3, diff_parts={K1 S2 I1 C0 F0}, diff_note={concepts: aromaticity and acidity of squaric acid | steps: squaric acid dianion has equal C-C bonds and 2 pi electrons so it is aromatic; squaric acid is stronger than carbonic acid; it releases CO2 from bicarbonate; check the carbonic acid statement}]
Incorrect statement among following is:

\begin{qoptions}
\opt{A}{All C---C bond in squaric acid dianion is of same bond length}
\opt{B}{Conjugated base of squaric acid is aromatic}
\opt{C}{Squaric acid on treatment with NaHCO$_3$ give effervescence of CO$_2$}
\opt{D}{Carbonic acid is a stronger acid than squaric acid}
\end{qoptions}
\end{bankq}
\begin{bankq}[id=SII19-P1-CHEM-S1-Q15, source={Score-II - Test 19 (20 Feb 2024)}, paper=1, qno=15, type=single, subtopic={Electronic effects}, answer={A}, answer_src=key, difficulty=2, diff_parts={K1 S1 I1 C0 F0}, diff_note={concepts: most stable resonance structure | steps: octet complete, fewer charges, negative charge on the more electronegative atom}]
Among the following most stable resonating structure is :-
\begin{qoptions}
\opt{A}{\bfigi{45.9}{SII19_p1_chem_s1_q15_1.png}}
\opt{B}{\bfigi{40.9}{SII19_p1_chem_s1_q15_2.png}}
\opt{C}{\bfigi{39.2}{SII19_p1_chem_s1_q15_5.png}}
\opt{D}{\bfigi{40.9}{SII19_p1_chem_s1_q15_3.png}}
\end{qoptions}
\end{bankq}
\begin{bankq}[id=SII19-P1-CHEM-S1-Q20, source={Score-II - Test 19 (20 Feb 2024)}, paper=1, qno=20, type=single, subtopic={Acidity and basicity}, answer={B}, answer_src=key, difficulty=3, diff_parts={K1 S2 I1 C0 F0}, diff_note={concepts: incorrect statement on halides and reactivity | steps: benzyl chloride undergoes SN1 and SN2; alpha halo ketone reactivity; check}]
Which of the following statement is not correct ?
\begin{qoptions}
\opt{A}{Compound Ph$-$CH$_2-$Cl can undergo S$_\mathrm N$1 \& S$_\mathrm N$2 reactions}
\opt{B}{For compound $\mathrm{R-\overset{\overset{O}{\|}}{C}-CH_2-Cl}$ \ Rate of S$_\mathrm N$2 $<$ Rate of S$_\mathrm N$1}
\opt{C}{Rate of $<$ S$_\mathrm N$1 \bfigi{57.0}{SII19_p1_chem_s1_q20_2.png}}
\opt{D}{Compound CH$_3$ $-$Cl will prefer to undergo S$_\mathrm N$2 when reacting with aq. KOH}
\end{qoptions}
\end{bankq}
\begin{bankq}[id=SII19-P1-CHEM-S2-Q10, source={Score-II - Test 19 (20 Feb 2024)}, paper=1, qno=10, type=numerical, subtopic={Electronic effects}, answer={4}, answer_src=key, difficulty=4, diff_parts={K1 S2 I1 C0 F1}, diff_note={concepts: number of groups showing -M or -R effect | steps: groups with pi bonds to electronegative atoms are -M; donors are +M; alkyl +I; count}]
How many species is/are $-$M/$-$R groups ?\\
(a) $-$OH\quad (b) $-$NO$_2$\quad (c) $-$CN\\
(d) $-$CH$_3$\quad (e) $-$O$^{\ominus}$\quad (f) $-\overset{\oplus}{\mathrm N}$H$_3$\\
(g) \bfigi{30.9}{SII19_p1_chem_s2_q10_1.png}\quad (h) $-$CH=NH
\end{bankq}
\begin{bankq}[id=SII22-P1-CHEM-S2-Q08, source={Score-II - Test 22 (14 Mar 2024)}, paper=1, qno=8, type=numerical, subtopic={Electronic effects}, answer={6}, answer_src=key, difficulty=3, diff_parts={K1 S1 I1 C0 F1}, diff_note={concepts: aromaticity | steps: check planarity cyclic conjugation and 4n+2 pi electrons for each compound}]
How many of the following compounds is are aromatic ?
\bfig{124.3}{SII22_p1_chem_s2_q08_1.png}
\end{bankq}
\begin{bankq}[id=SII25-P1-CHEM-S2-Q03, source={Score-II - Test 25 (20 Mar 2024)}, paper=1, qno=3, type=numerical, subtopic={Electronic effects}, answer={Bonus}, answer_src=solved, working={Printed key omits, truncates or corrupts this answer and the question was not solved; banked as Bonus.}, difficulty=4, diff_parts={K1 S1 I1 C1 F1}, diff_note={concepts: count aromatic compounds | steps: check planarity, cyclic conjugation and 4n + 2 for each}]
Find number of compound which are aromatic in nature :
\bfig{159.2}{SII25_p1_chem_s2_q03_1.png}
\end{bankq}
\begin{bankq}[id=SII26-P1-CHEM-S2-Q06, source={Score-II - Test 26 (20 Mar 2024)}, paper=1, qno=6, type=multi, subtopic={Electronic effects}, answer={A,B,C,D}, answer_src=key, difficulty=2, diff_parts={K1 S1 I0 C0 F1}, diff_note={concepts: correctness of drawn organic transformation options | steps: identify the major product by the usual mechanism for each of four reactions; compare with the drawn option}, also_in={TMAS - Test 01 (09 Jul 2023) P2 Q2}]
Correct option amongst the following is/are :
\begin{qoptions}
\opt{A}{\bfigi{109.4}{SII26_p1_chem_s2_q06_A.png}}
\opt{B}{\bfigi{128.8}{SII26_p1_chem_s2_q06_B.png}}
\opt{C}{\bfigi{160.9}{SII26_p1_chem_s2_q06_C.png}}
\opt{D}{\bfigi{209.1}{SII26_p1_chem_s2_q06_4.png}}
\end{qoptions}
\end{bankq}
\begin{bankq}[id=SII27-P1-CHEM-S1-Q12, source={Score-II - Test 27 (23 Mar 2024)}, paper=1, qno=12, type=single, subtopic={Acidity and basicity}, answer={D}, answer_src=key, difficulty=4, diff_parts={K1 S2 I1 C1 F0}, diff_note={concepts: decreasing order of acidic strength of drawn acids | steps: substituent effects on the conjugate base; rank A B C D; match option}]
\bfig{153.3}{SII27_p1_chem_s1_q12_1.png}
Which of the following is correct decreasing order of acidic strength ?
\begin{qoptions}
\opt{A}{B $>$ A $>$ D $>$ C}
\opt{B}{D $>$ A $>$ B $>$ C}
\opt{C}{A $>$ B $>$ C $>$ D}
\opt{D}{D $>$ B $>$ C $>$ A}
\end{qoptions}
\end{bankq}
\begin{bankq}[id=SII30-P1-CHEM-S1-Q14, source={Score-II - Test 30 (27 Mar 2024)}, paper=1, qno=14, type=single, subtopic={Acidity and basicity}, answer={D}, answer_src=key, difficulty=3, diff_parts={K1 S2 I1 C0 F0}, diff_note={concepts: strongest Bronsted base among drawn species | steps: basicity relates to stability of conjugate acid; amide anion strongest; rank}]
Which of the following is strongest Bronsted base?
\begin{qoptions}
\opt{A}{\bfigi{23.0}{SII30_p1_chem_s1_q14_1.png}}
\opt{B}{\bfigi{63.0}{SII30_p1_chem_s1_q14_2.png}}
\opt{C}{\bfigi{23.0}{SII30_p1_chem_s1_q14_3.png}}
\opt{D}{\bfigi{23.0}{SII30_p1_chem_s1_q14_4.png}}
\end{qoptions}
\end{bankq}
\begin{bankq}[id=SII31-P1-CHEM-S1-Q08, source={Score-II - Test 31 (29 Mar 2024)}, paper=1, qno=8, type=single, subtopic={Electronic effects}, answer={D}, answer_src=key, difficulty=3, diff_parts={K1 S1 I1 C1 F0}, diff_note={concepts: find the non aromatic species | steps: check planarity, conjugation and 4n + 2 rule for each; list non aromatic}]
Which of the following is/are non aromatic in nature ?
\bfig{145.8}{SII31_p1_chem_s1_q08_1.png}
\begin{qoptions}
\opt{A}{Only A}
\opt{B}{A and C}
\opt{C}{A, C and D}
\opt{D}{A, B, C and D}
\end{qoptions}
\end{bankq}
\begin{bankq}[id=SII33-P2-CHEM-S1-Q03, source={Score-II - Test 33 (01 Apr 2024)}, paper=2, qno=3, type=multi, subtopic={Electronic effects}, answer={A,B,D}, answer_src=key, difficulty=4, diff_parts={K1 S2 I1 C0 F1}, diff_note={concepts: resonance structures of methyl isocyanate | steps: charge separation structures of the N=C=O group; formal charges; valid resonance forms; check four}]
Which is/are the resonance structure of methyl isocyanate ?
\begin{qoptions}
\opt{A}{\bfigi{73.9}{SII33_p2_chem_s1_q03_A.png}}
\opt{B}{\bfigi{70.6}{SII33_p2_chem_s1_q03_B.png}}
\opt{C}{\bfigi{75.4}{SII33_p2_chem_s1_q03_C.png}}
\opt{D}{\bfigi{75.1}{SII33_p2_chem_s1_q03_D.png}}
\end{qoptions}
\end{bankq}
\begin{bankq}[id=SII38-P1-CHEM-S1-Q01, source={Score-II - Test 38 (12 Feb 2025)}, paper=1, qno=1, type=single, subtopic={Electronic effects}, answer={D}, answer_src=key, difficulty=4, diff_parts={K1 S2 I1 C1 F0}, diff_note={concepts: incorrect statement on organic reaction reagents and rates | steps: identify reagents P and Q from products; relative rate of nucleophilic addition; check the stated reagent in each of four options}]
Which of the following statement is \textbf{INCORRECT} ?
\begin{qoptions}
\opt{A}{\bfigi{34.4}{SII38_p1_chem_s1_q01_A.png} $\xrightarrow[\text{(ii) H}^+]{\text{(i) P}}$ \bfigi{41.6}{SII38_p1_chem_s1_q01_1.png}\\ In above reaction compound P is \bfigi{36.0}{SII38_p1_chem_s1_q01_7.png}}
\opt{B}{\bfigi{60.4}{SII38_p1_chem_s1_q01_B.png}\\ \bfigi{44.5}{SII38_p1_chem_s1_q01_4.png}\\ In above reaction compound Q = \bfigi{36.2}{SII38_p1_chem_s1_q01_8.png}}
\opt{C}{Order of rate of nucleophilic addition reaction is\\ \bfigi{115.4}{SII38_p1_chem_s1_q01_5.png}}
\opt{D}{IUPAC name of compound \bfigi{69.6}{SII38_p1_chem_s1_q01_6.png}\\ 3-Amino-2-formyl cyclopropane carboxylic acid}
\end{qoptions}
\end{bankq}
\begin{bankq}[id=SII41-P1-CHEM-S1-Q12, source={Score-II - Test 41 (19 Mar 2025)}, paper=1, qno=12, type=single, subtopic={Acidity and basicity}, answer={D}, answer_src=key, difficulty=2, diff_parts={K1 S1 I1 C0 F0}, diff_note={concepts: order of acidic strength of drawn acids | steps: substituents effect on conjugate base stability; rank}]
Give the correct order of Acidic strength.
\bfig{132.9}{SII41_p1_chem_s1_q12_1.png}
\begin{qoptions}
\opt{A}{(1) $>$ (2) $>$ (3) $>$ (4)}
\opt{B}{(3) $>$ (1) $>$ (2) $>$ (4)}
\opt{C}{(2) $>$ (1) $>$ (4) $>$ (3)}
\opt{D}{(2) $>$ (1) $>$ (3) $>$ (4)}
\end{qoptions}
\end{bankq}
\begin{bankq}[id=SII42-P1-CHEM-S2-Q05, source={Score-II - Test 42 (19 Mar 2025)}, paper=1, qno=5, type=numerical, subtopic={Electronic effects}, answer={10.00}, answer_src=key, difficulty=3, diff_parts={K1 S1 I1 C0 F1}, diff_note={concepts: resonance structures of a naphthol derivative | steps: draw the resonance forms by moving the lone pair of oxygen into the ring and the second ring; count the distinct structures}]
The number of resonating structures for organic compound P is

\begin{center}
\bfigi{135.4}{SII42_p2chem_n5_naphthol.png}
\end{center}
\end{bankq}
\begin{bankq}[id=SII45-P1-CHEM-S1-Q03, source={Score-II - Test 45 (26 Mar 2025)}, paper=1, qno=3, type=single, subtopic={Reactive intermediates}, answer={A}, answer_src=key, difficulty=2, diff_parts={K1 S1 I1 C0 F0}, diff_note={concepts: antiaromatic compound statement | steps: 4n pi electrons planar cyclic conjugated is antiaromatic; check the drawn compound}]
\textbf{Statement I :-} Compound \bfigi{25.2}{SII45_p1_chem_s1_q03_1.png} is antiaromatic in nature.\\
\textbf{Statement II :-} Compounds having (4n) $\pi$-electrons and planer with a cyclic conjugated $\pi$-system is anti aromatic.
\begin{qoptions}
\opt{A}{Statement I is false but Statement II is true}
\opt{B}{Both Statement I and Statement II are true}
\opt{C}{Both Statement I and Statement II are false}
\opt{D}{Statement I is true but Statement II is false}
\end{qoptions}
\end{bankq}
\begin{bankq}[id=SII46-P1-CHEM-S1-Q01, source={Score-II - Test 46 (26 Mar 2025)}, paper=1, qno=1, type=multi, subtopic={Electronic effects}, answer={A,C,D}, answer_src=key, difficulty=3, diff_parts={K1 S2 I0 C0 F1}, diff_note={concepts: reactions of phenols with hydroxylamine and tests on carbonyl compounds | steps: only phloroglucinol shows keto form with hydroxylamine; iodoform stoichiometry for a polymethyl ketone; check each statement and structure drawn}]
Which of the following statement(s) is/are \textbf{incorrect} ?
\begin{qoptions}
\opt{A}{\bfigi{276.3}{SII46_p1_chem_s1_q01_A.png}}
\opt{B}{1,3,5-trihydroxybenzene reacts with NH$_2$OH but 1,2,4-trihydroxybenzene and 1,2,3-trihydroxybenzene do not react with NH$_2$OH.}
\opt{C}{\bfigi{227.3}{SII46_p1_chem_s1_q01_C.png}}
\opt{D}{\bfigi{75.4}{SII46_p1_chem_s1_q01_D.png} number of Iodoform molecules produced per molecule of A is 3.}
\end{qoptions}
\end{bankq}
\begin{bankq}[id=SII47-P1-CHEM-S1-Q09, source={Score-II - Test 47 (30 Mar 2025)}, paper=1, qno=9, type=single, subtopic={Acidity and basicity}, answer={C}, answer_src=key, difficulty=2, diff_parts={K1 S1 I1 C0 F0}, diff_note={concepts: acidity of different hydrogens | steps: compare the stability of the conjugate bases by resonance and inductive effects}]
Most acidic hydrogen among the following is :
\bfig{186.7}{SII47_p1_chem_s1_q09_1.png}
\begin{qoptions}
\opt{A}{a}
\opt{B}{b}
\opt{C}{c}
\opt{D}{d}
\end{qoptions}
\end{bankq}

\section*{Score Advanced papers}
\begin{bankq}[id=SA04-P2-CHEM-S2-Q02, source={Score Advanced - Test 04 (04 Jul 2021)}, paper=2, qno=2, type=numerical, subtopic={Reactive intermediates}, answer={2}, answer_src=solved, working={Two 1,2-hydride shifts reach the tertiary cation}, difficulty=6, diff_parts={K2 S1 I2 C1 F1}, diff_note={concepts: number of shifts in a carbocation, reactive intermediates | steps: primary to secondary by H shift; secondary to tertiary by H shift}]
How many 1,2-shift(s) take(s) place in the following carbocation.\par
\bfig{65.3}{SA04_p2_chem_s2_q02_1.png}
\end{bankq}
\begin{bankq}[id=SA05-P1-CHEM-S1-Q19, source={Score Advanced - Test 05 (15 Jul 2021)}, paper=1, qno=19, type=single, subtopic={Acidity and basicity}, answer={B}, answer_src=key, difficulty=5, diff_parts={K2 S1 I2 C1 F0}, diff_note={concepts: incorrect statement on basicity of heterocycles, acidity and basicity | steps: check each statement on the structures}]
\bfig{150.6}{SA05_p1_chem_s1_q19_1.png}
Incorrect statement regarding above compounds is
\begin{qoptions}
\opt{A}{Basic strength I $>$ II $>$ III}
\opt{B}{Compound III is more basic than \bfigi{43.1}{SA05_p1_chem_s1_q19_B.png}}
\opt{C}{All are examples of heterocyclic aromatic compounds}
\opt{D}{Compound II shows inductive, hyperconjugation, resonance effect.}
\end{qoptions}
\end{bankq}
\begin{bankq}[id=SA07-P2-CHEM-S2-Q09, source={Score Advanced - Test 07 (16 Apr 2023)}, paper=2, qno=9, type=single, subtopic={Electronic effects}, answer={D}, answer_src=key, difficulty=4, diff_parts={K1 S2 I1 C1 F0}, diff_note={concepts: matching of reactions | steps: first reaction product type and its mechanism; second reaction; third; match each entry}]
\par
{\scriptsize\setlength{\tabcolsep}{1.5pt}\noindent\begin{tabular}{|L{.6cm}|L{3.1cm}|L{.6cm}|L{3.4cm}|}
\hline
 & List - I\par Order & & List - II\par Correct Parameter \\ \hline
(P) & \bfigi{66.4}{SA07_p2_chem_s2_q09_1.png} & (1) & Reactivity towards an electrophile or Lewis acid. \\ \hline
(Q) & \bfigi{79.4}{SA07_p2_chem_s2_q09_2.png} & (2) & Resonance Energy \\ \hline
(R) & \bfigi{68.6}{SA07_p2_chem_s2_q09_3.png} & (3) & At least one of the compound is aromatic \\ \hline
(S) & \bfigi{63.0}{SA07_p2_chem_s2_q09_4.png} & (4) & Enthalpy of hydrogenation \\ \hline
\end{tabular}}\par
\begin{qoptions}
\opt{A}{P $\rightarrow$ 2,3;Q $\rightarrow$ 3;R $\rightarrow$ 1,2,3;S $\rightarrow$ 1,2,3}
\opt{B}{P $\rightarrow$ 3,4;Q $\rightarrow$ 1;R $\rightarrow$ 1,4;S $\rightarrow$ 2,3,4}
\opt{C}{P $\rightarrow$ 2;Q $\rightarrow$ 1,2;R $\rightarrow$ 3,4;S $\rightarrow$ 2}
\opt{D}{P $\rightarrow$ 1,4;Q $\rightarrow$ 3;R $\rightarrow$ 2,3;S $\rightarrow$ 2,3}
\end{qoptions}
\end{bankq}
\begin{bankq}[id=SA08-P1-CHEM-S1-Q02, source={Score Advanced - Test 08 (22 Apr 2023)}, paper=1, qno=2, type=single, subtopic={Electronic effects}, answer={B}, answer_src=solved, working={Option B is wrong: pyrrole is more reactive than thiophene}, difficulty=1, diff_parts={K1 S1 I0 C0 F0}, diff_note={concepts: reactivity order | steps: compare the stability of the intermediates; check each order}]
Which of the following is not correct order of reactivity.
\begin{qoptions}
\opt{A}{\bfigi{133.7}{SA08_p1_chem_s1_q02_A.png}}
\opt{B}{\bfigi{133.0}{SA08_p1_chem_s1_q02_B.png}}
\opt{C}{\bfigi{127.5}{SA08_p1_chem_s1_q02_C.png}}
\opt{D}{\bfigi{135.4}{SA08_p1_chem_s1_q02_D.png}}
\end{qoptions}
\end{bankq}
\begin{bankq}[id=SA08-P1-CHEM-S3-Q06, source={Score Advanced - Test 08 (22 Apr 2023)}, paper=1, qno=6, type=numerical, subtopic={Acidity and basicity}, answer={13}, answer_src=key, difficulty=3, diff_parts={K1 S2 I0 C0 F1}, diff_note={concepts: basic strength comparison | steps: compare each compound's basicity by electronic effects; rank; serial number}]
Consider the following compound :\par
\bfigi{40.4}{SA08_p1_chem_s3_q06_2.png}\bfigi{40.2}{SA08_p1_chem_s3_q06_4.png}\bfigi{57.3}{SA08_p1_chem_s3_q06_3.png}\bfigi{38.5}{SA08_p1_chem_s3_q06_1.png}\par
\bfigi{40.6}{SA08_p1_chem_s3_q06_7.png}\bfigi{55.9}{SA08_p1_chem_s3_q06_5.png}\bfigi{38.4}{SA08_p1_chem_s3_q06_6.png}\bfigi{46.8}{SA08_p1_chem_s3_q06_8.png}\par
Serial number of most basic compound is (P)\par
Serial number of most acidic compound is (Q)\par
Total number of aromatic compound = (R)\par
Number of compound is represented in bracket.\par
Find the value of P + Q + R.
\end{bankq}
\begin{bankq}[id=SA08-P2-CHEM-S2-Q06, source={Score Advanced - Test 08 (22 Apr 2023)}, paper=2, qno=6, type=numerical, subtopic={Reactive intermediates}, answer={12}, answer_src=solved, working={Major product 2,3-dimethylbut-2-ene has 12 alpha hydrogens}, difficulty=4, diff_parts={K1 S2 I1 C0 F1}, diff_note={concepts: alpha hydrogens in the product | steps: identify the product; locate alpha carbons; count hydrogens}]
Identify total number of $\alpha$-`H' on major product of following reaction?\par
\bfig{156.5}{SA08_p2_chem_s2_q06_1.png}
\end{bankq}
\begin{bankq}[id=SA09-P2-CHEM-S1-Q04, source={Score Advanced - Test 09 (29 Apr 2023)}, paper=2, qno=4, type=multi, subtopic={Electronic effects}, answer={A,C,D}, answer_src=key, difficulty=5, diff_parts={K1 S2 I1 C1 F1}, diff_note={concepts: incorrect orders of nucleophilicity acidity and related properties | steps: nucleophilicity order in protic solvent; acid strength comparisons; stability orders; check each}]
Incorrect option among the following is/are
\begin{qoptions}
\opt{A}{Nucleophilicity order $\Rightarrow$ PhS$^\ominus$ $>$ CH$_3^\ominus$ $>$ I$^\ominus$ $>$ EtO$^\ominus$}
\opt{B}{Acidic strength order $\Rightarrow$ CH$_3$CO$_2$H $>$ H$_2$CO$_3$ $>$ \bfigi{115.1}{SA09_p2_chem_s1_q04_1.png}}
\opt{C}{Basic strength order $\Rightarrow$ \bfigi{108.0}{SA09_p2_chem_s1_q04_2.png}}
\opt{D}{Stability order : $\overset{+}{\text{C}}\text{Cl}_3$ $>$ $\overset{+}{\text{C}}\text{F}_3$ $>$ $\overset{+}{\text{C}}\text{H}_3$ $>$ CMe$_3$}
\end{qoptions}
\end{bankq}
\begin{bankq}[id=SA11-P1-CHEM-S1-Q04, source={Score Advanced - Test 11 (09 May 2023)}, paper=1, qno=4, type=multi, subtopic={Electronic effects}, answer={A,B}, answer_src=key, difficulty=2, diff_parts={K1 S1 I0 C0 F1}, diff_note={concepts: resonance energy comparison | steps: aromatic heterocycles versus butadiene; compare}, also_in={Score Advanced - Test 29 (01 May 2024) P1 Q4}]
Which of the following molecule(s) have higher resonance energy than 1,3$-$butadiene?
\begin{qoptions}
\opt{A}{Pyrrole}
\opt{B}{Thiophene}
\opt{C}{Tetrahydro Furan}
\opt{D}{1,4$-$cyclohexadiene}
\end{qoptions}
\end{bankq}
\begin{bankq}[id=SA11-P1-CHEM-S2-Q07, source={Score Advanced - Test 11 (09 May 2023)}, paper=1, qno=7, type=single, subtopic={Acidity and basicity}, answer={A}, answer_src=key, difficulty=1, diff_parts={K1 S1 I0 C0 F0}, diff_note={concepts: properties of bicarbonates | steps: structure of KHCO3 NaHCO3 LiHCO3; anion type}]
Match the property of the compound given in List$-$II with the compound given in List$-$I.\par
{\footnotesize\setlength{\tabcolsep}{1.5pt}\noindent\begin{tabular}{|c|l|c|L{4.3cm}|}
\hline
\multicolumn{2}{|c|}{List $-$ I} & \multicolumn{2}{c|}{List$-$ II}\\ \hline
(A) & KHCO$_3$ & (P) & contain polymeric anion\\ \hline
(B) & NaHCO$_3$ & (Q) & contain dimeric anion\\ \hline
(C) & LiHCO$_3$ & (R) & contain hydrogen bond\\ \hline
(D) & Mg(HCO$_3$)$_2$ & (S) & liberates CO$_2$ on heating\\ \hline
 & & (T) & can be prepared in solid state\\ \hline
\end{tabular}}\par
\begin{qoptions}
\opt{A}{A $\longrightarrow$ QRST;B $\longrightarrow$ PRST;C $\longrightarrow$ S;D $\longrightarrow$ ST}
\opt{B}{A $\longrightarrow$ QRST;B $\longrightarrow$ PST;C $\longrightarrow$ PS;D $\longrightarrow$ RST}
\opt{C}{A $\longrightarrow$ QST;B $\longrightarrow$ PST;C $\longrightarrow$ S;D $\longrightarrow$ RST}
\opt{D}{A $\longrightarrow$ QST;B $\longrightarrow$ PRST;C $\longrightarrow$ PS;D $\longrightarrow$ ST}
\end{qoptions}
\end{bankq}
\begin{bankq}[id=SA12-P1-CHEM-S1-Q02, source={Score Advanced - Test 12 (11 May 2023)}, paper=1, qno=2, type=single, subtopic={Acidity and basicity}, answer={C}, answer_src=key, difficulty=1, diff_parts={K1 S1 I0 C0 F0}, diff_note={concepts: statements about compounds and tests | steps: check each compound's test result; choose}, also_in={Score Advanced - Test 29 (01 May 2024) P2 Q2}]
Consider the following compounds and decide as to which of the following statements are true?\par
\bfig{271.3}{SA12_p1_chem_s1_q02_1.png}

\begin{qoptions}
\opt{A}{(II) gives no reaction with Na metal, however, 1 mole of (IV) on reaction with Na metal will liberate 22.4 litres of H$_2$ gas at 1atm \& 273K}
\opt{B}{(I) will give brisk effervescence on addition of NaHCO$_3$ but will not bring any change in the colour of Br$_2$ water}
\opt{C}{(III) liberates H$_2$ gas with Na metal, gives white precipitate with Tollen's reagent but does not respond towards Lucas reagent or 2, 4-DNP test.}
\opt{D}{(IV) gives turbidity with anhydrous ZnCl$_2$ + HCl}
\end{qoptions}
\end{bankq}
\begin{bankq}[id=SA15-P1-CHEM-S3-Q04, source={Score Advanced - Test 15 (18 May 2023)}, paper=1, qno=4, type=numerical, subtopic={Electronic effects}, answer={6}, answer_src=key, difficulty=5, diff_parts={K1 S2 I1 C1 F1}, diff_note={concepts: stability comparison of structure pairs | steps: resonance hyperconjugation and aromaticity; compare each pair; count}]
In how many of the below options, the first structure is more stable than the second.\par
\bfig{289.0}{SA15_p1_chem_s3_q04_1.png}
\end{bankq}
\begin{bankq}[id=SA18-P1-CHEM-S1-Q04, source={Score Advanced - Test 18 (27 May 2023)}, paper=1, qno=4, type=multi, subtopic={Reactive intermediates}, answer={A,B,C,D}, answer_src=key, difficulty=4, diff_parts={K1 S1 I1 C1 F1}, diff_note={concepts: order of stability of organic species | steps: carbocation or radical or alkene stability by hyperconjugation and resonance; compare each pair}]
Which of the following order of stability is correct?
\begin{qoptions}
\opt{A}{\bfigi{147.0}{SA18_p1_chem_s1_q04_A.png}}
\opt{B}{\bfigi{166.6}{SA18_p1_chem_s1_q04_B.png}}
\opt{C}{\bfigi{187.4}{SA18_p1_chem_s1_q04_C.png}}
\opt{D}{\bfigi{149.9}{SA18_p1_chem_s1_q04_D.png}}
\end{qoptions}
\end{bankq}
\begin{bankq}[id=SA18-P2-CHEM-S3-Q05, source={Score Advanced - Test 18 (27 May 2023)}, paper=2, qno=5, type=numerical, subtopic={Electronic effects}, answer={5}, answer_src=key, difficulty=4, diff_parts={K1 S1 I1 C1 F1}, diff_note={concepts: count aromatic species | steps: apply Huckel to each cyclic conjugated planar species including ions; count}]
How many of of the following species are aromatic\par
\bfig{236.3}{SA18_p2_chem_s3_q05_1.png}
\bfig{203.3}{SA18_p2_chem_s3_q05_2.png}
\end{bankq}
\begin{bankq}[id=SA22-P1-CHEM-S1-Q03, source={Score Advanced - Test 22 (10 Apr 2024)}, paper=1, qno=3, type=single, subtopic={Reactive intermediates}, answer={D}, answer_src=key, difficulty=4, diff_parts={K1 S2 I1 C1 F0}, diff_note={concepts: major product of the drawn organic reaction | steps: identify reagents; apply regiochemistry and stereochemistry; choose option}]
\bfig{121.7}{SA22_p1_chem_s1_q03_1.png}
\begin{qoptions}
\opt{A}{\bfigi{121.3}{SA22_p1_chem_s1_q03_A.png}}
\opt{B}{\bfigi{245.2}{SA22_p1_chem_s1_q03_B.png}}
\opt{C}{\bfigi{159.7}{SA22_p1_chem_s1_q03_C.png}}
\opt{D}{\bfigi{152.6}{SA22_p1_chem_s1_q03_D.png}}
\end{qoptions}
\end{bankq}
\begin{bankq}[id=SA24-P2-CHEM-S2-Q08, source={Score Advanced - Test 24 (17 Apr 2024)}, paper=2, qno=8, type=multi, subtopic={Electronic effects}, answer={A,D}, answer_src=key, difficulty=6, diff_parts={K1 S3 I1 C1 F1}, diff_note={concepts: molecules with all pi orbitals in the same plane | steps: benzene pi system is perpendicular to the plane; allene has perpendicular pi systems; cumulenes with odd numbers of double bonds; conjugated diene planar; check each}]
Select molecules(s) in which all the $\pi$-orbitals are in same plane.
\begin{qoptions}
\opt{A}{C$_6$H$_6$}
\opt{B}{O = C = C = C = O}
\opt{C}{H$_2$C = C = CH$_2$}
\opt{D}{CH$_2$ = CH $-$ CH = CH$_2$}
\end{qoptions}
\end{bankq}
\begin{bankq}[id=SA32-P1-CHEM-S3-Q04, source={Score Advanced - Test 32 (08 May 2024)}, paper=1, qno=4, type=numerical, subtopic={Electronic effects}, answer={6}, answer_src=key, difficulty=5, diff_parts={K1 S2 I1 C1 F1}, diff_note={concepts: stability comparison of structure pairs | steps: resonance hyperconjugation and aromaticity; compare each pair; count}]
In how many of the below options, the first structure is more stable than the second.\par
\bfig{114.2}{SA32_p1_chem_s3_q04_1.png}
\bfig{148.4}{SA32_p1_chem_s3_q04_2.png}
\bfig{121.0}{SA32_p1_chem_s3_q04_3.png}
\bfig{159.4}{SA32_p1_chem_s3_q04_4.png}
\end{bankq}

\section*{SRGBT papers}
\begin{bankq}[id=SRGBT07-P1-CHEM-S1-Q21, source={SRGBT - Test 07 (05 Sep 2021)}, paper=1, qno=21, type=multi, subtopic={Electronic effects}, answer={A}, answer_src=key, difficulty=3, diff_parts={K1 S1 I1 C0 F1}, diff_note={concepts: bond orders in anthracene from its resonance structures | steps: count the Kekule structures with a double bond at each position; bond 4 has the largest double-bond character}]
Which of the following C --C bond length is minimum in anthracene?\par
\bfig{75.2}{SRGBT07_p1_chem_s0_q21_1.png}
\begin{qoptions}
\opt{A}{4}
\opt{B}{3}
\opt{C}{2}
\opt{D}{1}
\end{qoptions}
\end{bankq}
\begin{bankq}[id=SRGBT07-P1-CHEM-S1-Q22, source={SRGBT - Test 07 (05 Sep 2021)}, paper=1, qno=22, type=multi, subtopic={Electronic effects}, answer={C}, answer_src=key, difficulty=2, diff_parts={K1 S1 I0 C0 F1}, diff_note={concepts: electronic effects on the compared property | steps: compare each pair in the given orders by resonance and inductive effects; identify the order that is reversed}]
Which of the following order is incorrect ?\par
\bfig{245.2}{SRGBT07_p1_chem_s0_q22_1.png}
\end{bankq}
\begin{bankq}[id=SRGBT07-P1-CHEM-S1-Q26, source={SRGBT - Test 07 (05 Sep 2021)}, paper=1, qno=26, type=multi, subtopic={Acidity and basicity}, answer={B}, answer_src=key, difficulty=2, diff_parts={K1 S1 I0 C0 F1}, diff_note={concepts: acid-base reactions go towards the weaker acid | steps: propyne (pKa 25) cannot be deprotonated by hydroxide since water (pKa 15.7) is the stronger acid}]
Which of the following acid base reaction is not fesible.
\begin{qoptions}
\opt{A}{\bfigi{200.3}{SRGBT07_p1_chem_s0_q26_A.png}}
\opt{B}{Me $-$ C $\equiv$ CH + NaOH $\to$ MeC $\equiv\overset{-}{\mathrm C}\overset{\oplus}{\mathrm{Na}}$ + H$_2$O}
\opt{C}{\bfigi{182.7}{SRGBT07_p1_chem_s0_q26_2.png}}
\opt{D}{Me$_3$CONa + H$_2$O $\to$ Me$_3$COH + NaOH}
\end{qoptions}
\end{bankq}
\begin{bankq}[id=SRGBT07-P1-CHEM-S1-Q27, source={SRGBT - Test 07 (05 Sep 2021)}, paper=1, qno=27, type=multi, subtopic={Acidity and basicity}, answer={A}, answer_src=key, difficulty=3, diff_parts={K1 S1 I1 C0 F1}, diff_note={concepts: availability of the nitrogen lone pair (delocalisation into a ring or carbonyl) | steps: compare lone-pair delocalisation in the structures; the one with the lone pair in an aromatic sextet is least basic}]
Which one of the following is least basic in character ?\par
\bfig{293.5}{SRGBT07_p1_chem_s0_q27_1.png}
\end{bankq}
\begin{bankq}[id=SRGBT07-P1-CHEM-S1-Q28, source={SRGBT - Test 07 (05 Sep 2021)}, paper=1, qno=28, type=multi, subtopic={Reactive intermediates}, answer={A}, answer_src=key, difficulty=3, diff_parts={K1 S1 I1 C0 F1}, diff_note={concepts: comparison of the shown structures by resonance and inductive effects | steps: rank the structures by their electronic effects; pick option A}]
\bfig{328.0}{SRGBT07_p1_chem_s0_q28_1.png}
\end{bankq}
\begin{bankq}[id=SRGBT07-P1-CHEM-S1-Q29, source={SRGBT - Test 07 (05 Sep 2021)}, paper=1, qno=29, type=multi, subtopic={Reactive intermediates}, answer={A,D}, answer_src=key, difficulty=5, diff_parts={K2 S2 I1 C0 F1}, diff_note={concepts: carbocation and radical stability by hyperconjugation and resonance, aromaticity | steps: count hyperconjugative hydrogens and resonance for each species; check aromatic stabilisation; mark A and D correct}]
Which of the given following stability order is/are correct.\par
\bfig{316.8}{SRGBT07_p1_chem_s0_q29_1.png}
\end{bankq}
\begin{bankq}[id=SRGBT07-P1-CHEM-S1-Q30, source={SRGBT - Test 07 (05 Sep 2021)}, paper=1, qno=30, type=multi, subtopic={Acidity and basicity}, answer={B,C}, answer_src=key, difficulty=5, diff_parts={K2 S2 I1 C0 F1}, diff_note={concepts: acidity from conjugate base stability (resonance and -I effects), ortho effect | steps: stabilise each conjugate base; compare -M and -I substituents; check each proposed order}]
Correct increasing order of acidic strength of following compounds is / are ?\par
\bfig{245.2}{SRGBT07_p1_chem_s0_q30_1.png}
\end{bankq}
\begin{bankpassage}[id=SRGBT07-P1-CHEM-S1-G35-36, source={SRGBT - Test 07 (05 Sep 2021)}, questions={35-36}]
Consider the following figure representing variation of heat of combustion of isomeric butenes.\par
\bfig{197.9}{SRGBT07_p1_chem_s0_q34_3.png}
Based on this answer the following questions.
\end{bankpassage}
\begin{bankq}[id=SRGBT07-P1-CHEM-S1-Q35, source={SRGBT - Test 07 (05 Sep 2021)}, paper=1, qno=35, type=multi, subtopic={Electronic effects}, answer={D}, answer_src=key, difficulty=2, diff_parts={K1 S1 I0 C0 F1}, diff_note={concepts: vector addition of bond dipoles in isomers | steps: the symmetric isomer S has dipoles cancelling so minimum}, group=SRGBT07-P1-CHEM-S1-G35-36]
Minimum dipole moment is of isomers.
\begin{qoptions}
\opt{A}{P}
\opt{B}{Q}
\opt{C}{R}
\opt{D}{S}
\end{qoptions}
\end{bankq}
\begin{bankq}[id=SRGBT07-P1-CHEM-S1-Q36, source={SRGBT - Test 07 (05 Sep 2021)}, paper=1, qno=36, type=multi, subtopic={Electronic effects}, answer={B}, answer_src=key, difficulty=3, diff_parts={K1 S1 I1 C0 F1}, diff_note={concepts: rate of HBr addition set by carbocation stability | steps: Q forms the most stabilised carbocation so it reacts fastest}, group=SRGBT07-P1-CHEM-S1-G35-36]
Which gives fastest reaction with HBr?
\begin{qoptions}
\opt{A}{P}
\opt{B}{Q}
\opt{C}{R}
\opt{D}{S}
\end{qoptions}
\end{bankq}
\begin{bankq}[id=SRGBT07-P1-CHEM-S1-Q38, source={SRGBT - Test 07 (05 Sep 2021)}, paper=1, qno=38, type=multi, subtopic={Electronic effects}, answer={10}, answer_src=key, difficulty=4, diff_parts={K1 S2 I0 C1 F1}, diff_note={concepts: drawing resonance structures by moving the charge through conjugation | steps: move the charge over each conjugated position; include structures on the substituent; count 10}]
The number of resoance structures for Z is:\par
\bfig{117.2}{SRGBT07_p1_chem_s0_q38_1.png}
\end{bankq}
\begin{bankq}[id=SRGBT07-P1-CHEM-S1-Q39, source={SRGBT - Test 07 (05 Sep 2021)}, paper=1, qno=39, type=multi, subtopic={Reactive intermediates}, answer={3}, answer_src=key, difficulty=3, diff_parts={K1 S1 I1 C0 F1}, diff_note={concepts: keto-enol tautomerism exchanges alpha hydrogens with deuterium | steps: count the enolisable alpha hydrogens including those through conjugation; 3}]
The number of deuterium atoms found in the tautomer of \bfigi{44.2}{SRGBT07_p1_chem_s0_q38_2.png} when it is kept in OD$^\ominus$ / D$_2$O for a long period of time?
\end{bankq}
\begin{bankq}[id=SRGBT07-P1-CHEM-S1-Q40, source={SRGBT - Test 07 (05 Sep 2021)}, paper=1, qno=40, type=multi, subtopic={Electronic effects}, answer={541}, answer_src=key, difficulty=4, diff_parts={K2 S2 I0 C0 F1}, diff_note={concepts: mesomeric effects of substituent groups, hyperconjugation of alkyl groups | steps: -M groups NO NO2 OCOCH3 amide carbonyl and SR2+ count 5; +M groups NO OCOCH3 OH and SH count 4; +H only for ethyl; answer 541}]
Consider the following groups :\par
$-$NO, $-\overset{+}{\mathrm S}\mathrm R_2$, $-$NO$_2$, $-$CMe$_3$, $-$OCOCH$_3$, $-\mathrm C(=\mathrm O)-$NHCH$_3$, $-\overset{+}{\mathrm N}$H$_3$, $-\overset{\ominus}{\mathrm B}$F$_3$, $-$CH$_2$--CH$_3$, $-$OH, $-$SH\par
Among given compounds.\par
(a) Number of groups which can show $-$M effect = P\par
(b) Number of groups which can show +M effect = Q\par
(c) Number of group which can show +H effect = R\par
\textbf{Give your answer in 3 digits as PQR.}
\end{bankq}
\begin{bankq}[id=SRGBT14-P1-CHEM-S1-Q21, source={SRGBT - Test 14 (24 Oct 2021)}, paper=1, qno=21, type=multi, subtopic={Electronic effects}, answer={C}, answer_src=key, difficulty=4, diff_parts={K2 S2 I0 C0 F1}, diff_note={concepts: stability and dipole moment of the given species, heats of combustion and aromatic bond orders | steps: check the stability and dipole orders; heat of combustion order of the isomers is reversed; cyclopentadienyl anion bond order 1.4 is correct}]
Select incorrect option.
\begin{qoptions}
\opt{A}{\bfigi{101.0}{SRGBT14_p1_chem_s0_q21_A.png} is correct order of stability.}
\opt{B}{\bfigi{94.0}{SRGBT14_p1_chem_s0_q21_B.png} is correct order of dipole moment.}
\opt{C}{\bfigi{195.0}{SRGBT14_p1_chem_s0_q21_C.png} is correct order of heat of combustion.}
\opt{D}{Bond order of C--C bond in conjugate base of cyclo penta--1,3--diene is 7/5.}
\end{qoptions}
\end{bankq}
\begin{bankq}[id=SRGBT14-P1-CHEM-S1-Q22, source={SRGBT - Test 14 (24 Oct 2021)}, paper=1, qno=22, type=multi, subtopic={Acidity and basicity}, answer={D}, answer_src=key, difficulty=2, diff_parts={K1 S1 I0 C0 F1}, diff_note={concepts: carboxylic acid strength versus phenols and sulphonic acids | steps: the phenol-type compound D is weaker than benzoic acid}]
Which of the following is less acidic than benzoic acid.
\begin{qoptions}
\opt{A}{HCOOH}
\opt{B}{PhSO$_3$H}
\opt{C}{\bfigi{50.2}{SRGBT14_p1_chem_s0_q22_C.png}}
\opt{D}{\bfigi{51.7}{SRGBT14_p1_chem_s0_q22_D.png}}
\end{qoptions}
\end{bankq}
\begin{bankq}[id=SRGBT14-P1-CHEM-S1-Q28, source={SRGBT - Test 14 (24 Oct 2021)}, paper=1, qno=28, type=multi, subtopic={Acidity and basicity}, answer={C}, answer_src=key, difficulty=2, diff_parts={K1 S1 I0 C0 F1}, diff_note={concepts: acidity from conjugate base stabilisation | steps: compare resonance and inductive stabilisation of the four anions; order ii > iii > i > iv}]
(i) \bfigi{35.0}{SRGBT14_p1_chem_s0_q28_i.png} (ii) \bfigi{35.3}{SRGBT14_p1_chem_s0_q28_ii.png} (iii) \bfigi{44.9}{SRGBT14_p1_chem_s0_q28_iii.png} (iv) \bfigi{47.6}{SRGBT14_p1_chem_s0_q28_iv.png}\par
Correct order of acidic nature is:
\begin{qoptions}
\opt{A}{iii $>$ i $>$ ii $>$ iv}
\opt{B}{i $>$ iii $>$ ii $>$ iv}
\opt{C}{ii $>$ iii $>$ i $>$ iv}
\opt{D}{ii $>$ iii $>$ iv $>$ i}
\end{qoptions}
\end{bankq}
\begin{bankq}[id=SRGBT14-P1-CHEM-S1-Q30, source={SRGBT - Test 14 (24 Oct 2021)}, paper=1, qno=30, type=multi, subtopic={Electronic effects}, answer={A,B,C,D}, answer_src=key, difficulty=2, diff_parts={K1 S1 I0 C0 F1}, diff_note={concepts: enol content from conjugation hydrogen bonding and aromaticity | steps: compare the keto and enol stabilisation in each pair; all four comparisons hold}]
Which of the following is/are correct for enol content ?
\begin{qoptions}
\opt{A}{\bfigi{69.1}{SRGBT14_p1_chem_s0_q30_A.png}}
\opt{B}{\bfigi{115.3}{SRGBT14_p1_chem_s0_q30_B.png}}
\opt{C}{\bfigi{104.8}{SRGBT14_p1_chem_s0_q30_C.png}}
\opt{D}{\bfigi{51.8}{SRGBT14_p1_chem_s0_q30_D.png}}
\end{qoptions}
\end{bankq}
\begin{bankpassage}[id=SRGBT14-P1-CHEM-S1-G35-36, source={SRGBT - Test 14 (24 Oct 2021)}, questions={35-36}]
In organic compounds electrons of $\sigma$ \& $\pi$--bonds are show displacement from their normal position due to various reason. It is known as electronic displacement effect. For electronic displacement effect give answer of following questions.
\end{bankpassage}
\begin{bankq}[id=SRGBT14-P1-CHEM-S1-Q35, source={SRGBT - Test 14 (24 Oct 2021)}, paper=1, qno=35, type=multi, subtopic={Electronic effects}, answer={D}, answer_src=key, difficulty=2, diff_parts={K1 S1 I0 C0 F1}, diff_note={concepts: inductive effect depends on distance and geometry | steps: the +I magnitude depends on geometry so C is incorrect}, group=SRGBT14-P1-CHEM-S1-G35-36]
For inductive effect, select incorrect statement.
\begin{qoptions}
\opt{A}{\bfigi{98.8}{SRGBT14_p1_chem_s0_q35_A.png}, inductive effect of $-$Me on $-$COOH is same in both compounds.}
\opt{B}{Inductive effect is distance dependent effect.}
\opt{C}{Magnitude of +I effect does not depend upon geometry of compound.}
\opt{D}{Maximum cationic group can show $-$I effect but it is not necessary for all.}
\end{qoptions}
\end{bankq}
\begin{bankq}[id=SRGBT14-P1-CHEM-S1-Q36, source={SRGBT - Test 14 (24 Oct 2021)}, paper=1, qno=36, type=multi, subtopic={Electronic effects}, answer={D}, answer_src=key, difficulty=2, diff_parts={K1 S1 I0 C0 F1}, diff_note={concepts: mesomeric effects of groups | steps: some groups such as -NO can show both +M and -M so D}, group=SRGBT14-P1-CHEM-S1-G35-36]
Select correct statement for mesomeric effect.
\begin{qoptions}
\opt{A}{All catonic groups can show $-$M effect.}
\opt{B}{All catonic groups can show $+$M effect.}
\opt{C}{$\dot{\mathrm C}\mathrm H_2-\mathrm{CN}$ is less stable than $\dot{\mathrm C}\mathrm H_3$ due to $-$M effect of $-$CN group.}
\opt{D}{Some groups can show both $+$M and $-$M depending upon condition.}
\end{qoptions}
\end{bankq}
\begin{bankq}[id=SRGBT14-P1-CHEM-S1-Q37, source={SRGBT - Test 14 (24 Oct 2021)}, paper=1, qno=37, type=multi, subtopic={Electronic effects}, answer={11}, answer_src=key, difficulty=4, diff_parts={K1 S2 I0 C1 F1}, diff_note={concepts: counting resonance structures by moving the charge through conjugation | steps: delocalise the charge over the rings; include structures on the substituent; total 11}]
How many resonating structures are possible for \bfigi{50.4}{SRGBT14_p1_chem_s0_q37_1.png} including this structure ?
\end{bankq}
\begin{bankq}[id=SRGBT14-P1-CHEM-S1-Q40, source={SRGBT - Test 14 (24 Oct 2021)}, paper=1, qno=40, type=multi, subtopic={Acidity and basicity}, answer={13}, answer_src=key, difficulty=4, diff_parts={K2 S2 I0 C0 F1}, diff_note={concepts: basicity and acidity of the listed compounds, aromaticity counting | steps: most basic compound P and most acidic Q from lone-pair availability and anion stability; count aromatic compounds R; P + Q + R = 13}]
Consider the following compound :\par
\bfig{294.7}{SRGBT14_p1_chem_s0_q40_1.png}

Serial number of most basic compound is (P)\par
Serial number of most acidic compound is (Q)\par
Total number of aromatic compound = (R)\par
Number of compound is represented in bracket.\par
Find the value of P + Q + R.
\end{bankq}

\section*{Nurture SRG papers}
\begin{bankq}[id=NSRG03-P1-CHEM-S1-Q01, source={Nurture SRG - Test 03 (04 Feb 2024)}, paper=1, qno=1, type=single, subtopic={Acidity and basicity}, answer={C}, answer_src=key, difficulty=2, diff_parts={K1 S2 I0 C0 F0}, diff_note={concepts: acidity of hydrocarbons | steps: stability of the conjugate bases of each; aromatic anion; compare to the alkene}]
Which of the following is INCORRECT statement ?
\begin{qoptions}
\opt{A}{\bfigi{164.9}{NSRG03_p1_chem_s1_q01_A.png}}
\opt{B}{\bfigi{187.6}{NSRG03_p1_chem_s1_q01_B.png}}
\opt{C}{\bfigi{203.8}{NSRG03_p1_chem_s1_q01_C.png}}
\opt{D}{\bfigi{35.1}{NSRG03_p1_chem_s1_q01_D.png}\par is more acidic than $\mathrm{CH_2{=}CH{-}CH_2{-}CH{=}CH_2}$}
\end{qoptions}
\end{bankq}
\begin{bankq}[id=NSRG03-P1-CHEM-S1-Q03, source={Nurture SRG - Test 03 (04 Feb 2024)}, paper=1, qno=3, type=single, subtopic={Reactive intermediates}, answer={B}, answer_src=key, difficulty=3, diff_parts={K1 S2 I1 C0 F0}, diff_note={concepts: reaction intermediates | steps: identify reagent; the intermediate formed; stability comparison}]
During the following reaction, which intermediate is dominantly appeared ?\par
\bfig{182.2}{NSRG03_p1_chem_s1_q03_1.png}
\begin{qoptions}
\opt{A}{\bfigi{114.2}{NSRG03_p1_chem_s1_q03_A.png}}
\opt{B}{\bfigi{88.2}{NSRG03_p1_chem_s1_q03_B.png}}
\opt{C}{\bfigi{105.6}{NSRG03_p1_chem_s1_q03_C.png}}
\opt{D}{\bfigi{102.1}{NSRG03_p1_chem_s1_q03_D.png}}
\end{qoptions}
\end{bankq}
\begin{bankq}[id=NSRG03-P1-CHEM-S3-Q05, source={Nurture SRG - Test 03 (04 Feb 2024)}, paper=1, qno=5, type=numerical, subtopic={Acidity and basicity}, answer={6}, answer_src=solved, working={Soluble in NaOH: o-toluic acid, p-cresol, 2,6-dihydroxybenzoic acid, picric acid, 4-hydroxypyridine and the cyanonaphthol}, difficulty=2, diff_parts={K1 S1 I0 C0 F1}, diff_note={concepts: acidity and solubility in NaOH | steps: identify phenols and acids among compounds; count soluble}]
How many compounds are soluble in aqueous NaOH ?\par
\bfig{207.5}{NSRG03_p1_chem_s3_q05_1.png}
\end{bankq}
\begin{bankq}[id=NSRG03-P1-CHEM-S3-Q06, source={Nurture SRG - Test 03 (04 Feb 2024)}, paper=1, qno=6, type=numerical, subtopic={Electronic effects}, answer={5}, answer_src=solved, working={Aromatic: cyclopropenyl cation, tropylium cation, cyclopentadienyl anion, 1,2-dihydronaphthalene (benzene ring) and phenanthrene}, difficulty=2, diff_parts={K1 S1 I0 C0 F1}, diff_note={concepts: aromaticity | steps: check each ring for conjugation planarity and 4n+2; count}]
Among the following, the number of aromatic compound(s) is(are)\par
\bfig{193.0}{NSRG03_p1_chem_s3_q06_1.png}
\end{bankq}
\begin{bankq}[id=NSRG03-P2-CHEM-S1-Q02, source={Nurture SRG - Test 03 (04 Feb 2024)}, paper=2, qno=2, type=single, subtopic={Reactive intermediates}, answer={D}, answer_src=solved, working={The cation ring expands to an oxygen stabilised cation; the methyl shift is not the major mode}, difficulty=3, diff_parts={K1 S2 I1 C0 F0}, diff_note={concepts: carbocation rearrangement statements | steps: identify the cation; ring expansion; stability gain; check each statement}]
The INCORRECT statement Regarding Rearrangement in the given carbocation?\par
\bfig{68.5}{NSRG03_p2_chem_s1_q02_1.png}
\begin{qoptions}
\opt{A}{It rearranges to most stable carbocation by ring Expansion}
\opt{B}{Its most stable form obtained after rearrangement has a resonance structure with completed valence shells of all the atoms.}
\opt{C}{To its most stable form obtained after rearrangement, oxygen decreases stability by inductive effect but increases stability by resonance.}
\opt{D}{$\mathrm{CH_3}$-shift is the major mode of Rearrangement.}
\end{qoptions}
\end{bankq}
\begin{bankq}[id=NSRG03-P2-CHEM-S1-Q03, source={Nurture SRG - Test 03 (04 Feb 2024)}, paper=2, qno=3, type=single, subtopic={Electronic effects}, answer={B}, answer_src=solved, working={Enolisation puts sp2 carbons in a four membered ring (strained), so the ordinary cyclohexanone has the highest enol content, then the saturated cyclobutanone, then the unsaturated one}, difficulty=1, diff_parts={K1 S1 I0 C0 F0}, diff_note={concepts: enol content | steps: compare the stability of the enol forms by conjugation and hydrogen bonding; order}]
Among the following compounds;\par
\bfig{156.0}{NSRG03_p2_chem_s1_q03_1.png}
The correct decreasing order of Enol-content is;
\begin{qoptions}
\opt{A}{I, II, III}
\opt{B}{III, I, II}
\opt{C}{II, I, III}
\opt{D}{II, III, I}
\end{qoptions}
\end{bankq}
\begin{bankq}[id=NSRG03-P2-CHEM-S2-Q05, source={Nurture SRG - Test 03 (04 Feb 2024)}, paper=2, qno=5, type=multi, subtopic={Acidity and basicity}, answer={A,B,C}, answer_src=solved, working={Proton sponge with donating OMe is more basic; ortho effect; picric acid > formic > thiophenol > phenol; the primary ammonium is the stronger acid so its pKa is lower, making the last statement false}, difficulty=2, diff_parts={K1 S1 I0 C0 F1}, diff_note={concepts: order of properties | steps: check each of the four orders}]
Select the option which contains the CORRECT order against the parameter mentioned in the parenthesis
\begin{qoptions}
\opt{A}{\bfigi{154.6}{NSRG03_p2_chem_s2_q05_A.png}}
\opt{B}{\bfigi{206.3}{NSRG03_p2_chem_s2_q05_B.png}}
\opt{C}{\bfigi{209.1}{NSRG03_p2_chem_s2_q05_C.png}}
\opt{D}{\bfigi{208.9}{NSRG03_p2_chem_s2_q05_D.png}}
\end{qoptions}
\end{bankq}
\begin{bankq}[id=NSRG03-P2-CHEM-S2-Q10, source={Nurture SRG - Test 03 (04 Feb 2024)}, paper=2, qno=10, type=multi, subtopic={Electronic effects}, answer={A,B}, answer_src=solved, working={Cyclopropenone has a large dipole (aromatic cyclopropenium oxide); acetone is both electrophile and nucleophile; the isopropylidene compound has no geometrical isomers; alkoxide is more basic than thiolate}, difficulty=3, diff_parts={K1 S2 I0 C0 F1}, diff_note={concepts: reactivity statements | steps: acetone as electrophile and nucleophile; alkoxide basicity; check each}]
Which of these statement are correct?
\begin{qoptions}
\opt{A}{\bfigi{196.2}{NSRG03_p2_chem_s2_q10_A.png}}
\opt{B}{Acetone behaves as electrophile as well as nucleophile}
\opt{C}{\bfigi{185.2}{NSRG03_p2_chem_s2_q10_C.png}}
\opt{D}{$\mathrm{RO^{\ominus}}$ is less basic than $\mathrm{RS^{\ominus}}$}
\end{qoptions}
\end{bankq}

\section*{Enthusiast Pre-Score papers}
\begin{bankq}[id=EPS01-P1-CHEM-S1-Q04, source={Enthusiast Pre-Score - Test 01 (04 Nov 2020)}, paper=1, qno=4, type=single, subtopic={Electronic effects}, answer={A}, answer_src=key, difficulty=3, diff_parts={K1 S2 I1 C0 F0}, diff_note={concepts: dipoles in pyrrole and pyridine | steps: lone pair in the ring for pyrrole is delocalised; pyridine lone pair in sp2; direction of each dipole}]
The correct orientation of dipoles in pyrrole and pyridine is :\par
\begin{qoptions}
\opt{A}{\bfigi{105.4}{EPS01_p1_chem_s1_q04_A.png}}
\opt{B}{\bfigi{105.4}{EPS01_p1_chem_s1_q04_B.png}}
\opt{C}{\bfigi{105.4}{EPS01_p1_chem_s1_q04_C.png}}
\opt{D}{\bfigi{105.4}{EPS01_p1_chem_s1_q04_D.png}}
\end{qoptions}
\end{bankq}
\begin{bankq}[id=EPS01-P2-CHEM-S3-Q04, source={Enthusiast Pre-Score - Test 01 (04 Nov 2020)}, paper=2, qno=4, type=single, subtopic={Electronic effects}, answer={4}, answer_src=key, difficulty=1, diff_parts={K1 S1 I0 C0 F0}, diff_note={concepts: aromaticity | steps: check planarity and 4n+2 for each structure; count}]
Among the following, the number of aromatic compound(s) is \underline{\hspace{0.8cm}}\par
\bfig{360.2}{EPS01_p2_chem_s3_q04_1.png}
\end{bankq}
\begin{bankq}[id=EPS03-P2-CHEM-S1-Q03, source={Enthusiast Pre-Score - Test 03 (23 Nov 2020)}, paper=2, qno=3, type=multi, subtopic={Electronic effects}, answer={A,C}, answer_src=key, difficulty=2, diff_parts={K1 S1 I0 C0 F1}, diff_note={concepts: resonance structure stability | steps: charge separation and octet in each pair; pick pairs with the more stable first}]
Among the following resonating structures, in which option(s), (I) is more stable than (II)\par
\bfig{371.8}{EPS03_p2_chem_s1_q03_B.png}
\end{bankq}

\section*{JEE Advanced (2021 Series) papers}
\begin{bankq}[id=JA02-P1-CHEM-S2-Q08, source={JEE Advanced (2021 Series) - Test 02 (04 Feb 2021)}, paper=1, qno=8, type=multi, subtopic={Acidity and basicity}, answer={A,B,C}, answer_src=key, difficulty=6, diff_parts={K1 S3 I1 C1 F1}, diff_note={concepts: acidity conjugate base properties of compounds I to V | steps: delocalisation stabilises conjugate bases; aromatic conjugate base of cyclopentadiene type; nitro substituent increases acidity; order of acidity; check each statement}]
With respect to the compounds \textbf{I-V}, choose the correct statement(s).
\bfig{364.2}{JA02_p1_chem_s2_q08_1.png}

\begin{qoptions}
\opt{A}{The acidity of compound \textbf{I} is due to delocalization in the conjugate base.}
\opt{B}{The conjugate base of compound \textbf{IV} is aromatic.}
\opt{C}{Compound \textbf{II} becomes more acidic, when it has a -NO$_2$ substituent.}
\opt{D}{The acidity of compounds follows the order \textbf{I $>$ IV $>$ V $>$ II $>$ III.}}
\end{qoptions}
\end{bankq}
\begin{bankq}[id=JA02-P1-CHEM-S3-Q06, source={JEE Advanced (2021 Series) - Test 02 (04 Feb 2021)}, paper=1, qno=6, type=numerical, subtopic={Reactive intermediates}, answer={18.00}, answer_src=key, difficulty=6, diff_parts={K1 S2 I2 C1 F1}, diff_note={concepts: degree unsaturation of a cyclised product in sulphuric acid | steps: ester of triphenyl benzoic acid forms acylium; intramolecular Friedel Crafts attack onto ortho phenyl gives fluorenone type coloured compound; count rings and pi bonds; DBE | I: recognise intramolecular cyclisation to a fluorenone}]
In the following reaction, compound \textbf{Q} is obtained from compound \textbf{P} via an ionic intermediate
\bfig{307.6}{JA02_p1_chem_s3_q06_1.png}

What is the degree of unsaturation of Q ?
\end{bankq}
\begin{bankq}[id=JA02-P2-CHEM-S1-Q05, source={JEE Advanced (2021 Series) - Test 02 (04 Feb 2021)}, paper=2, qno=5, type=multi, subtopic={Acidity and basicity}, answer={C,D}, answer_src=key, difficulty=5, diff_parts={K2 S2 I1 C0 F1}, diff_note={concepts: basicity of amines with resonance, steric and inductive effects | steps: compare lone pair availability in I to IV; resonance delocalisation reduces basicity in III relative to IV where steric inhibition of resonance exists; order of basicity; pKb differences}]
Consider the following four compounds \textbf{I}, \textbf{II}, \textbf{III}, and \textbf{IV}.
\begin{center}
\bfigi{31.0}{JA02_p2_chem_s1_q05_4.png}\hspace{0.5em}\bfigi{69.4}{JA02_p2_chem_s1_q05_3.png}\hspace{0.5em}\bfigi{83.5}{JA02_p2_chem_s1_q05_2.png}\hspace{0.5em}\bfigi{83.4}{JA02_p2_chem_s1_q05_1.png}
\end{center}
Choose the correct statement(s).
\begin{qoptions}
\opt{A}{The order of basicity is \textbf{II} $>$ \textbf{I} $>$ \textbf{III} $>$ \textbf{IV}.}
\opt{B}{The magnitude of p$K_\mathrm{b}$ difference between \textbf{I} and \textbf{II} is more than that between \textbf{III} and \textbf{IV}.}
\opt{C}{Resonance effect is more in \textbf{III} than in \textbf{IV}.}
\opt{D}{Steric effect makes compound \textbf{IV} more basic than \textbf{III}.}
\end{qoptions}
\end{bankq}
\begin{bankq}[id=JA03-P1-CHEM-S3-Q04, source={JEE Advanced (2021 Series) - Test 03 (01 Mar 2021)}, paper=1, qno=4, type=single, subtopic={Acidity and basicity}, answer={5.00}, answer_src=key, difficulty=5, diff_parts={K1 S3 I1 C1 F0}, diff_note={concepts: acidity comparison with carbonic acid | steps: pKa of carbonic acid near 6.4; squaric acid is strongly acidic; carboxylic acids and phenols with electron withdrawing groups compare; propyne and methanol are weaker; H2S pKa 7 is weaker; count those stronger}]
How many of the following species/compounds are more acidic than carbonic acid?
\begin{center}
(a) Squaric acid \bfigi{80.9}{JA03_p1_chem_s3_q04_1.png}\par\medskip
(b) \bfigi{34.3}{JA03_p1_chem_s3_q04_2.png}\quad (c) \bfigi{55.0}{JA03_p1_chem_s3_q04_3.png}\par\medskip
(d) \bfigi{62.6}{JA03_p1_chem_s3_q04_4.png}\quad (e) \bfigi{80.3}{JA03_p1_chem_s3_q04_5.png}\par\medskip
(f) \bfigi{30.4}{JA03_p1_chem_s3_q04_6.png}\quad (g) HC$\equiv$C$-$CH$_3$\par\medskip
(h) MeOH\quad (i) H$_2$S
\end{center}
\end{bankq}
\begin{bankq}[id=JA04-P1-CHEM-S1-Q03, source={JEE Advanced (2021 Series) - Test 04 (04 Apr 2021)}, paper=1, qno=3, type=single, subtopic={Electronic effects}, answer={D}, answer_src=key, difficulty=3, diff_parts={K1 S2 I1 C0 F0}, diff_note={concepts: stability of canonical structures by charge separation and octet | steps: compare octet completion in each contributor; fewer charges and negative charge on more electronegative atom is more stable; rank all four}]
\begin{center}
(I)~\bfigi{31.6}{JA04_p1_chem_s1_q02_1.png}\quad (II)~\bfigi{35.5}{JA04_p1_chem_s1_q02_2.png}\par\medskip
(III)~\bfigi{31.6}{JA04_p1_chem_s1_q03_1.png}\quad (IV)~\bfigi{26.4}{JA04_p1_chem_s1_q03_2.png}
\end{center}
The correct order of stability for the given canonical structure is
\begin{qoptions}
\opt{A}{IV $>$ II $>$ III $>$ I}
\opt{B}{IV $>$ II $>$ I $>$ III}
\opt{C}{II $>$ IV $>$ III $>$ I}
\opt{D}{II $>$ IV $>$ I $>$ III}
\end{qoptions}
\end{bankq}
\begin{bankq}[id=JA06-P2-CHEM-S1-Q03, source={JEE Advanced (2021 Series) - Test 06 (22 Apr 2021)}, paper=2, qno=3, type=multi, subtopic={Acidity and basicity}, answer={B,C}, answer_src=key, difficulty=4, diff_parts={K1 S2 I1 C0 F1}, diff_note={concepts: acid strength comparison with the reference acid | steps: soluble in NaHCO3 means stronger than carbonic acid; compare with the given reference acid by substituent effects; carboxylic acids with electron withdrawing groups are stronger; pick those stronger than the reference}]
Which of the following will be soluble in $\mathrm{NaHCO_3}$ solution and more acidic than \bfigi{73.3}{JA06_p2_chem_s1_q03_1.png}
\begin{qoptions}
\opt{A}{\bfigi{33.2}{JA06_p2_chem_s1_q03_A.png}}
\opt{B}{\bfigi{91.6}{JA06_p2_chem_s1_q03_B.png}}
\opt{C}{\bfigi{41.6}{JA06_p2_chem_s1_q03_C.png}}
\opt{D}{\bfigi{44.9}{JA06_p2_chem_s1_q03_D.png}}
\end{qoptions}
\end{bankq}
\begin{bankq}[id=JA06-P2-CHEM-S3-Q05, source={JEE Advanced (2021 Series) - Test 06 (22 Apr 2021)}, paper=2, qno=5, type=single, subtopic={Electronic effects}, answer={10.00}, answer_src=key, difficulty=4, diff_parts={K1 S2 I1 C1 F0}, diff_note={concepts: delocalised and localised nitrogen lone pairs | steps: a nitrogen lone pair is delocalised if it is conjugated with a pi system or adjacent p orbital; check each nitrogen in the drawn structure; add the tag numbers of those delocalised}]
Sum of tags of Nitrogen lone pairs which are delocaized.
\bfig{111.5}{JA06_p2_chem_s3_q05_1.png}
\end{bankq}
\begin{bankq}[id=JA08-P1-CHEM-S1-Q07, source={JEE Advanced (2021 Series) - Test 08 (01 May 2021)}, paper=1, qno=7, type=multi, subtopic={Acidity and basicity}, answer={A,B,C,D}, answer_src=solved, working={salicylic > m- > p-hydroxybenzoic acid; 2-nitrophenol < 2,4-dinitrophenol; HF < HCl < HBr < HI; HOCl < HClO2 < HClO3 < HClO4.}, difficulty=5, diff_parts={K1 S2 I1 C1 F1}, diff_note={concepts: acid strength orders | steps: ortho effect makes salicylic acid strongest and para hydroxy weakest; nitro groups increase phenol acidity; hydrogen halide acidity increases down the group; oxyacid strength increases with oxidation state of chlorine}]
Identify correct order(s) of acidic strength?
\begin{qoptions}
\opt{A}{\bfigi{98.2}{JA08_p1_chem_s1_q07_A.png}}
\opt{B}{\bfigi{91.3}{JA08_p1_chem_s1_q07_B.png}}
\opt{C}{HF $<$ HCl $<$ HBr $<$ HI}
\opt{D}{HOCl $<$ HClO$_2$ $<$ HClO$_3$ $<$ HClO$_4$}
\end{qoptions}
\end{bankq}
\begin{bankq}[id=JA09-P2-CHEM-S1-Q02, source={JEE Advanced (2021 Series) - Test 09 (07 May 2021)}, paper=2, qno=2, type=multi, subtopic={Acidity and basicity}, answer={A,B,C}, answer_src=key, difficulty=5, diff_parts={K1 S3 I0 C1 F1}, diff_note={concepts: acid strength and conjugate base stability decide equilibrium direction | steps: reaction A ethoxide takes a proton from a more acidic C-H or O-H; reaction B proton goes from a stronger to a weaker acid by the structure shown; reaction C carbanion takes a proton from the amine giving the more stable amide; reaction D proton moves to the stronger base so it is backward; list pKa for each pair and keep the forward ones | C: four reactions each needing a pKa comparison from drawn structures}]
Among the given reaction which of the following favours forward reaction.
\begin{qoptions}
\opt{A}{EtO$^{\ominus}$ $+$ \bfigi{172.2}{JA09_p2_chem_s1_q02_1.png}}
\opt{B}{\bfigi{148.8}{JA09_p2_chem_s1_q02_B.png}}
\opt{C}{CH$_3-$CH$_2^{\ominus}$ $+$ CH$_3-$NH$_2$ $\rightleftharpoons$ CH$_3-$CH$_3$ $+$ CH$_3-$NH$^{\ominus}$}
\opt{D}{\bfigi{223.0}{JA09_p2_chem_s1_q02_D.png}}
\end{qoptions}
\end{bankq}
\begin{bankq}[id=JA12-P1-CHEM-S1-Q05, source={JEE Advanced (2021 Series) - Test 12 (25 May 2021)}, paper=1, qno=5, type=multi, subtopic={Reactive intermediates}, answer={A,B,C,D}, answer_src=solved, working={Homoallyl cation can be formed and shifts to the allylic, cyclopropylmethyl and cyclobutyl cations.}, difficulty=6, diff_parts={K1 S3 I2 C0 F1}, diff_note={concepts: formation and rearrangement of carbocations from homoallyl systems | steps: protonation of the alcohol and loss of water gives the primary homoallyl cation; this shifts a hydride to give the secondary allylic cation; it also closes to the cyclopropylmethyl cation; this rearranges to the cyclobutyl cation; all four are accessible | I: the homoallyl cyclopropylmethyl cyclobutyl cation manifold is a special non classical rearrangement}]
When \bfigi{53.4}{JA12_p1_chem_s1_q05_0.png} is treated with HCl then different carbocations formed can be
\begin{qoptions}
\opt{A}{\bfigi{35.5}{JA12_p1_chem_s1_q05_A.png}}
\opt{B}{\bfigi{38.2}{JA12_p1_chem_s1_q05_B.png}}
\opt{C}{\bfigi{43.0}{JA12_p1_chem_s1_q05_C.png}}
\opt{D}{\bfigi{24.5}{JA12_p1_chem_s1_q05_D.png}}
\end{qoptions}
\end{bankq}
\begin{bankq}[id=JA13-P2-CHEM-S2-Q05, source={JEE Advanced (2021 Series) - Test 13 (31 May 2021)}, paper=2, qno=5, type=multi, subtopic={Electronic effects}, answer={A,B,D}, answer_src=key, working={Printed key.}, difficulty=3, diff_parts={K1 S2 I0 C0 F1}, diff_note={concepts: bond lengths and dipole moments with resonance in vinyl chloride | steps: C-Cl in vinyl chloride is shorter because of resonance and sp2 carbon; C=C in vinyl chloride is longer because of delocalisation; dipole moment of ethyl chloride is larger than vinyl chloride}]
Which of the given statement is correct :
\begin{qoptions}
\opt{A}{The C--Cl bond length in $\mathrm{CH_3CH_2Cl}$ is more than that in and CH=CH-Cl}
\opt{B}{The C $=$ C bond length in $\mathrm{CH_2=CH_2}$ is less than that in and $\mathrm{CH_2=CHCl}$}
\opt{C}{The C $-$ C bond length in $\mathrm{CH_2=CH_2}$, $\mathrm{CH_3CH_2-Cl}$ and $\mathrm{CH_2=CHCl}$ is equal}
\opt{D}{The dipole moment of $\mathrm{CH_3CH_2Cl}$ is more than that in $\mathrm{CH_2=CHCl}$}
\end{qoptions}
\end{bankq}

\section*{TOAS papers}
\begin{bankq}[id=TOAS01-P1-CHEM-S2-Q08, source={TOAS - Test 01 (30 Jan 2022)}, paper=1, qno=8, type=multi, subtopic={Reactive intermediates}, answer={B,C,D}, answer_src=key, difficulty=3, diff_parts={K1 S2 I0 C0 F1}, diff_note={concepts: nucleophiles | steps: dichlorocarbene is an electrophile; amide ion is a nucleophile; alkene as pi nucleophile; last structure}]
Which of the following is/are nucleophile
\begin{qoptions}
\opt{A}{$:\mathrm{CCl_2}$}
\opt{B}{$\overset{\bullet\bullet}{\mathrm{N}}\mathrm{H_2}$}
\opt{C}{$\mathrm{CH_3{-}CH{=}CH_2}$}
\opt{D}{\bfigi{30.8}{TOAS01_p1_chem_s2_q08_D.png}}
\end{qoptions}
\end{bankq}
\begin{bankq}[id=TOAS01-P1-CHEM-S2-Q09, source={TOAS - Test 01 (30 Jan 2022)}, paper=1, qno=9, type=multi, subtopic={Reactive intermediates}, answer={A,D}, answer_src=key, difficulty=5, diff_parts={K2 S2 I1 C0 F1}, diff_note={concepts: aromaticity, reaction products | steps: product of each reaction; apply Huckel rule to each product; pick those that are aromatic | I: aromatic ion formed by hydride or proton loss}]
Which of the following reaction(s) gives aromatic product as major product :
\begin{qoptions}
\opt{A}{\bfigi{98.0}{TOAS01_p1_chem_s2_q09_A.png}}
\opt{B}{\bfigi{76.8}{TOAS01_p1_chem_s2_q09_B.png}}
\opt{C}{\bfigi{111.2}{TOAS01_p1_chem_s2_q09_C.png}}
\opt{D}{\bfigi{76.5}{TOAS01_p1_chem_s2_q09_D.png}}
\end{qoptions}
\end{bankq}
\begin{bankq}[id=TOAS01-P2-CHEM-S1-Q15, source={TOAS - Test 01 (30 Jan 2022)}, paper=2, qno=15, type=single, subtopic={Reactive intermediates}, answer={A}, answer_src=key, difficulty=3, diff_parts={K1 S2 I1 C0 F0}, diff_note={concepts: carbocation rearrangement | steps: first carbocation; successive 1,2 shifts; final cation; count shifts | I: shifts continue to the most stable cation}]
\bfigi{106.1}{TOAS01_p2_chem_s1_q15_1.png}\par
Identify total number of 1,2-shift for above reaction ?
\begin{qoptions}
\opt{A}{2}
\opt{B}{3}
\opt{C}{4}
\opt{D}{5}
\end{qoptions}
\end{bankq}
\begin{bankq}[id=TOAS01-P2-CHEM-S1-Q20, source={TOAS - Test 01 (30 Jan 2022)}, paper=2, qno=20, type=single, subtopic={Reactive intermediates}, answer={D}, answer_src=key, difficulty=3, diff_parts={K1 S2 I1 C0 F0}, diff_note={concepts: carbocation stability and rearrangement | steps: rearrange the cation by hydride or methyl shift; compare the final cations; choose the most stable | I: shift to give a tertiary or resonance stabilised cation}]
The most stable carbocation after rearrangement is
\bfig{51.2}{TOAS01_p2_chem_s1_q20_1.png}
\begin{qoptions}
\opt{A}{\bfigi{52.9}{TOAS01_p2_chem_s1_q20_A.png}}
\opt{B}{\bfigi{50.7}{TOAS01_p2_chem_s1_q20_B.png}}
\opt{C}{\bfigi{51.9}{TOAS01_p2_chem_s1_q20_C.png}}
\opt{D}{\bfigi{51.1}{TOAS01_p2_chem_s1_q20_D.png}}
\end{qoptions}
\end{bankq}
\begin{bankq}[id=TOAS01-P2-CHEM-S2-Q10, source={TOAS - Test 01 (30 Jan 2022)}, paper=2, qno=10, type=numerical, subtopic={Reactive intermediates}, answer={5}, answer_src=solved, working={Rearrangement possible (to a more stable cation): (i), (iii), (iv) by H shift, (vi) H shift to oxocarbenium, (vii) methyl shift to tertiary; benzyl, iminium-stabilised, cumyl, tertiary (ix) and vinyl do not}, difficulty=4, diff_parts={K1 S2 I1 C0 F1}, diff_note={concepts: carbocation rearrangement | steps: identify the cation type; adjacent hydrogen or alkyl available; count those that can rearrange | I: shift only if it gives a more stable cation}]
How many of the given carbocation can undergo rearrangement :
\bfig{150.1}{TOAS01_p2_chem_s2_q10_1.png}
\end{bankq}
\begin{bankq}[id=TOAS03-P1-CHEM-S2-Q07, source={TOAS - Test 03 (27 Feb 2022)}, paper=1, qno=7, type=multi, subtopic={Reactive intermediates}, answer={A,C,D}, answer_src=key, difficulty=3, diff_parts={K1 S2 I0 C0 F1}, diff_note={concepts: stability of intermediates | steps: intermediate in each option A and B; in C and D; compare stabilities}]
Which of following option is/are correct?
\begin{qoptions}
\opt{A}{\bfigi{110.8}{TOAS03_p1_chem_s2_q07_A.png}}
\opt{B}{\bfigi{173.6}{TOAS03_p1_chem_s2_q07_B.png}}
\opt{C}{\bfigi{179.2}{TOAS03_p1_chem_s2_q07_C.png}}
\opt{D}{\bfigi{180.0}{TOAS03_p1_chem_s2_q07_D.png}}
\end{qoptions}
\end{bankq}
\begin{bankq}[id=TOAS03-P2-CHEM-S2-Q01, source={TOAS - Test 03 (27 Feb 2022)}, paper=2, qno=1, type=numerical, subtopic={Reactive intermediates}, answer={6}, answer_src=solved, working={Rearrange: (b) cyclobutylmethyl (ring expansion), (e) isobutyl, (f) 2-hydroxycyclohexyl, (g) neopentyl, (h) secondary next to a tertiary C-H, (j) ring cation next to the isopropyl-bearing carbon; no rearrangement for (a) tertiary, (c) benzyl, (d) allylic, (i) oxocarbenium}, difficulty=3, diff_parts={K1 S2 I0 C0 F1}, diff_note={concepts: carbocation rearrangements | steps: classify each carbocation; check for a more stable rearranged cation; count 6}]
Total number of following carbocation which will undergo rearragement ?
\bfig{239.9}{TOAS03_p2_chem_s2_q01_1.png}
\end{bankq}
\begin{bankq}[id=TOAS03-P2-CHEM-S3-Q02, source={TOAS - Test 03 (27 Feb 2022)}, paper=2, qno=2, type=numerical, subtopic={Reactive intermediates}, answer={4}, answer_src=key, difficulty=3, diff_parts={K1 S2 I0 C0 F1}, diff_note={concepts: carbocation shifts | steps: first carbocation and shift; second shift; later shifts; count the 1,2 shifts}]
\bfigi{135.4}{TOAS03_p2_chem_s3_q02_1.png}

Total number of 1,2 shifts (rearrangements) involve in the above reaction.
\end{bankq}
\begin{bankq}[id=TOAS05-P2-CHEM-S1-Q03, source={TOAS - Test 05 (03 Apr 2022)}, paper=2, qno=3, type=single, subtopic={Reactive intermediates}, answer={B}, answer_src=key, difficulty=1, diff_parts={K1 S1 I0 C0 F0}, diff_note={concepts: carbocation intermediates | steps: identify the intermediate; rearranged product}]
\bfig{120.5}{TOAS05_p2_chem_s1_q03_1.png}
\begin{qoptions}
\opt{A}{\bfigi{25.0}{TOAS05_p2_chem_s1_q03_A.png}}
\opt{B}{\bfigi{25.2}{TOAS05_p2_chem_s1_q03_B.png}}
\opt{C}{\bfigi{25.5}{TOAS05_p2_chem_s1_q03_C.png}}
\opt{D}{$\mathrm{CH_2{=}CH_2}$}
\end{qoptions}
\end{bankq}
\begin{bankq}[id=TOAS05-P2-CHEM-S1-Q04, source={TOAS - Test 05 (03 Apr 2022)}, paper=2, qno=4, type=single, subtopic={Reactive intermediates}, answer={D}, answer_src=key, difficulty=3, diff_parts={K1 S2 I1 C0 F0}, diff_note={concepts: reaction sequence with rearrangement | steps: first step; rearrangement; final step | I: carbocation rearrangement in the sequence}]
\bfig{141.1}{TOAS05_p2_chem_s1_q04_1.png}
Final product ($B$) of the reaction sequence is :
\begin{qoptions}
\opt{A}{\bfigi{55.1}{TOAS05_p2_chem_s1_q04_A.png}}
\opt{B}{\bfigi{55.1}{TOAS05_p2_chem_s1_q04_B.png}}
\opt{C}{\bfigi{52.9}{TOAS05_p2_chem_s1_q04_C.png}}
\opt{D}{\bfigi{52.9}{TOAS05_p2_chem_s1_q04_D.png}}
\end{qoptions}
\end{bankq}
\begin{bankq}[id=TOAS06-P1-CHEM-S1-Q03, source={TOAS - Test 06 (17 Apr 2022)}, paper=1, qno=3, type=single, subtopic={Reactive intermediates}, answer={A}, answer_src=key, difficulty=1, diff_parts={K1 S1 I0 C0 F0}, diff_note={concepts: carbocation rearrangement | steps: carbocation formed; rearranged product}]
Identify the major product :
\bfig{93.0}{TOAS06_p1_chem_s1_q03_1.png}
\begin{qoptions}
\opt{A}{\bfigi{45.8}{TOAS06_p1_chem_s1_q03_A.png}}
\opt{B}{\bfigi{47.2}{TOAS06_p1_chem_s1_q03_B.png}}
\opt{C}{\bfigi{43.8}{TOAS06_p1_chem_s1_q03_C.png}}
\opt{D}{\bfigi{56.5}{TOAS06_p1_chem_s1_q03_D.png}}
\end{qoptions}
\end{bankq}
\begin{bankq}[id=TOAS06-P1-CHEM-S2-Q07, source={TOAS - Test 06 (17 Apr 2022)}, paper=1, qno=7, type=multi, subtopic={Electronic effects}, answer={B,D}, answer_src=key, difficulty=3, diff_parts={K1 S2 I0 C0 F1}, diff_note={concepts: electrophiles and nucleophilicity | steps: classify CO2, Na+, AlCl3; nucleophilicity order; options C and D}]
Which of following is/are correct :
\begin{qoptions}
\opt{A}{$\mathrm{CO_2}$, $\mathrm{Na^+}$, $\mathrm{AlCl_3}$ all are electrophiles.}
\opt{B}{$\mathrm{F^{\ominus}<Cl^{\ominus}<Br^{\ominus}<I^{\ominus}}$ : nucleophillicity order (in aqueous medium)}
\opt{C}{\bfigi{114.1}{TOAS06_p1_chem_s2_q07_C.png}}
\opt{D}{\bfigi{189.6}{TOAS06_p1_chem_s2_q07_D.png}}
\end{qoptions}
\end{bankq}
\begin{bankq}[id=TOAS06-P1-CHEM-S3-Q04, source={TOAS - Test 06 (17 Apr 2022)}, paper=1, qno=4, type=numerical, subtopic={Reactive intermediates}, answer={4}, answer_src=solved, working={Dehydration gives the carbocation; loss of the very acidic ring H gives 6-methylfulvene C$_7$H$_8$: DOU $=(2\cdot7+2-8)/2=4$}, difficulty=4, diff_parts={K1 S2 I1 C0 F1}, diff_note={concepts: dehydration and unsaturation | steps: carbocation formed; loss of H to give fulvene; DOU | I: aromatic cyclopentadienyl driving force}]
\bfigi{147.4}{TOAS06_p1_chem_s3_q04_1.png}\\
Degree of unsaturation in P
\end{bankq}
\begin{bankq}[id=TOAS11-P1-CHEM-S3-Q05, source={TOAS - Test 11 (27 Jun 2022)}, paper=1, qno=5, type=numerical, subtopic={Reactive intermediates}, answer={78.00}, answer_src=key, difficulty=4, diff_parts={K1 S2 I1 C0 F1}, diff_note={concepts: pinacol rearrangement | steps: X is the ring expanded ketone; LiAlH4 gives the alcohol Y; dehydration gives Z; molar mass | I: ring expansion of the pinacol}]
Observe the following reaction sequence
\bfig{304.7}{TOAS11_p1_chem_s3_q05_1.png}

Calculate molecular mass ($W$) of product Z and report your answer as $N$, where $N=\dfrac W3$
\end{bankq}
\begin{bankq}[id=TOAS13-P2-CHEM-S2-Q05, source={TOAS - Test 13 (24 Jul 2022)}, paper=2, qno=5, type=numerical, subtopic={Reactive intermediates}, answer={5}, answer_src=key, difficulty=3, diff_parts={K1 S2 I0 C0 F1}, diff_note={concepts: resonance in a carbocation | steps: delocalisation to the ortho and para ring carbons; extension through the C=C; count the new forms}]
How many resonating structures of carbocation `M' as shown below are possible (excluding given resonating structure) :
\bfig{53.2}{TOAS13_p2_chem_s2_q05_1.png}
\end{bankq}
\begin{bankq}[id=TOAS15-P2-CHEM-S1-Q03, source={TOAS - Test 15 (22 Aug 2022)}, paper=2, qno=3, type=multi, subtopic={Reactive intermediates}, answer={A,C}, answer_src=solved, working={Singlet $\mathrm{CF_2}$ is more stabilised than singlet $\mathrm{CCl_2}$ (2p-2p back donation); dimethoxycarbene singlet beats the triplet; for dimethylcarbene the triplet is the more stable form}, difficulty=4, diff_parts={K1 S2 I1 C0 F1}, diff_note={concepts: stability of singlet and triplet carbenes | steps: fluoro versus chloro carbene; alkyl carbene triplet; alkoxy carbene singlet}]
In which of the following I$^{\text{st}}$ is more stable than II$^{\text{nd}}$
\begin{qoptions}
\opt{A}{\bfigi{29.2}{TOAS15_p2_chem_s1_q03_A.png}\ \bfigi{32.0}{TOAS15_p2_chem_s1_q03_2.png}}
\opt{B}{\bfigi{87.2}{TOAS15_p2_chem_s1_q03_B.png}}
\opt{C}{\bfigi{84.6}{TOAS15_p2_chem_s1_q03_C.png}}
\opt{D}{None}
\end{qoptions}
\end{bankq}
\begin{bankq}[id=TOAS17-P2-CHEM-S1-Q06, source={TOAS - Test 17 (18 Sep 2022)}, paper=2, qno=6, type=single, subtopic={Reactive intermediates}, answer={C}, answer_src=key, difficulty=4, diff_parts={K2 S2 I1 C0 F0}, diff_note={concepts: reaction sequence with phenols, acylation | steps: first reaction with the reagent; second step; choose the product}]
\bfig{139.6}{TOAS17_p2_chem_s1_q06_1.png}
Major product of the reaction is:
\begin{qoptions}
\opt{A}{\bfigi{55.6}{TOAS17_p2_chem_s1_q06_A.png}}
\opt{B}{\bfigi{63.0}{TOAS17_p2_chem_s1_q06_B.png}}
\opt{C}{\bfigi{56.5}{TOAS17_p2_chem_s1_q06_C.png}}
\opt{D}{\bfigi{61.0}{TOAS17_p2_chem_s1_q06_D.png}}
\end{qoptions}
\end{bankq}
\begin{bankq}[id=TOAS19-P1-CHEM-S1-Q08, source={TOAS - Test 19 (05 Feb 2023)}, paper=1, qno=8, type=multi, subtopic={Acidity and basicity}, answer={A,B,C,D}, answer_src=key, difficulty=4, diff_parts={K2 S2 I0 C0 F1}, diff_note={concepts: nucleophilicity and aromatic cations, steric inhibition of resonance | steps: ion stability; basic strength with steric effects; carbocation rate of dehydration; halide nucleophilicity in protic solvents}]
Identify correct statement?
\begin{qoptions}
\opt{A}{F$^\ominus<$ Cl$^\ominus<$ Br$^\ominus<$ I$^\ominus$ Nucleophilicity order in H$_2$O}
\opt{B}{\bfigi{118.3}{TOAS19_p1_chem_s1_q08_B.png}}
\opt{C}{\bfigi{213.3}{TOAS19_p1_chem_s1_q08_C.png}}
\opt{D}{\bfigi{189.6}{TOAS19_p1_chem_s1_q08_D.png}}
\end{qoptions}
\end{bankq}

\section*{TOPS papers}
\begin{bankq}[id=TOPS03-P1-CHEM-S1-Q03, source={TOPS - Test 03 (20 Nov 2022)}, paper=1, qno=3, type=single, subtopic={Electronic effects}, answer={B}, answer_src=key, difficulty=3, diff_parts={K1 S2 I1 C0 F0}, diff_note={concepts: compound with inductive hyperconjugation and resonance effects | steps: identify the structural features for each compound; resonance requires a conjugated system; hyperconjugation requires alpha hydrogens; check four}]
Compound in which inductive, hyper conjugation and resonance all three effect are present?
\begin{qoptions}
\opt{A}{\bfigi{31.8}{TOPS03_p1_chem_s1_q03_A.png}}
\opt{B}{\bfigi{24.0}{TOPS03_p1_chem_s1_q03_B.png}}
\opt{C}{\bfigi{66.1}{TOPS03_p1_chem_s1_q03_C.png}}
\opt{D}{\bfigi{24.0}{TOPS03_p1_chem_s1_q03_D.png}}
\end{qoptions}
\end{bankq}
\begin{bankq}[id=TOPS03-P1-CHEM-S2-Q07, source={TOPS - Test 03 (20 Nov 2022)}, paper=1, qno=7, type=multi, subtopic={Acidity and basicity}, answer={A,B,C,D}, answer_src=key, difficulty=4, diff_parts={K1 S2 I1 C0 F1}, diff_note={concepts: acidic strength order of organic acids | steps: substituent effects inductive and resonance; stability of conjugate base; compare pairs}]
Identify correct acidic strength order?
\begin{qoptions}
\opt{A}{\bfigi{21.5}{TOPS03_p1_chem_s2_q07_1.png}$<$\bfigi{28.9}{TOPS03_p1_chem_s2_q07_A.png}}
\opt{B}{\bfigi{74.5}{TOPS03_p1_chem_s2_q07_B.png}}
\opt{C}{\bfigi{80.4}{TOPS03_p1_chem_s2_q07_C.png}}
\opt{D}{\bfigi{78.5}{TOPS03_p1_chem_s2_q07_D.png}}
\end{qoptions}
\end{bankq}
\begin{bankq}[id=TOPS03-P1-CHEM-S2-Q08, source={TOPS - Test 03 (20 Nov 2022)}, paper=1, qno=8, type=multi, subtopic={Acidity and basicity}, answer={A,B,C,D}, answer_src=key, difficulty=4, diff_parts={K1 S2 I1 C0 F1}, diff_note={concepts: basic strength order of amines | steps: inductive effects and resonance; hybridisation of nitrogen; compare pairs}]
Identify correct basic strength order :
\begin{qoptions}
\opt{A}{\bfigi{125.8}{TOPS03_p1_chem_s2_q08_A.png}}
\opt{B}{\bfigi{21.5}{TOPS03_p1_chem_s2_q08_2.png}$<$\bfigi{18.7}{TOPS03_p1_chem_s2_q08_B.png}}
\opt{C}{\bfigi{26.7}{TOPS03_p1_chem_s2_q08_4.png}$>$\bfigi{26.5}{TOPS03_p1_chem_s2_q08_C.png}}
\opt{D}{\bfigi{135.7}{TOPS03_p1_chem_s2_q08_D.png}}
\end{qoptions}
\end{bankq}
\begin{bankq}[id=TOPS03-P1-CHEM-S4-Q08, source={TOPS - Test 03 (20 Nov 2022)}, paper=1, qno=8, type=numerical, subtopic={Reactive intermediates}, answer={4}, answer_src=key, difficulty=3, diff_parts={K1 S1 I1 C0 F1}, diff_note={concepts: count tertiary hydrogens in a structure | steps: identify tertiary carbons; count hydrogens attached}]
Number of 3$^\circ$ hydrogens present in the given compound
\bfig{60.4}{TOPS03_p1_chem_s4_q08_1.png}
\end{bankq}
\begin{bankq}[id=TOPS03-P2-CHEM-S1-Q03, source={TOPS - Test 03 (20 Nov 2022)}, paper=2, qno=3, type=multi, subtopic={Reactive intermediates}, answer={B,D}, answer_src=key, difficulty=4, diff_parts={K1 S2 I1 C0 F1}, diff_note={concepts: resonance of positive charge or empty orbital | steps: identify carbocations with adjacent double bond or lone pair; resonance only if the vacant orbital is conjugated; check four}]
In which of the following ions, positive charge / vacant orbital is in resonance?
\begin{qoptions}
\opt{A}{\bfigi{26.7}{TOPS03_p2_chem_s1_q03_A.png}}
\opt{B}{\bfigi{24.0}{TOPS03_p2_chem_s1_q03_B.png}}
\opt{C}{\bfigi{26.7}{TOPS03_p2_chem_s1_q03_C.png}}
\opt{D}{\bfigi{24.0}{TOPS03_p2_chem_s1_q03_D.png}}
\end{qoptions}
\end{bankq}
\begin{bankq}[id=TOPS03-P2-CHEM-S1-Q04, source={TOPS - Test 03 (20 Nov 2022)}, paper=2, qno=4, type=multi, subtopic={Electronic effects}, answer={B,C}, answer_src=key, difficulty=5, diff_parts={K1 S3 I1 C0 F1}, diff_note={concepts: degree of unsaturation and hydrogens in cubane | steps: C8H8 gives DU 5; cubane's carbons are all tertiary with one H; count of tertiary hydrogens 8; no secondary carbons; check four}]
Statements which are correct about cubane
\begin{qoptions}
\opt{A}{D.U. $=6$}
\opt{B}{D.U. $=5$}
\opt{C}{No. of 3$^\circ$ H $=8$}
\opt{D}{No. of 2$^\circ$C $=4$}
\end{qoptions}
\end{bankq}
\begin{bankq}[id=TOPS03-P2-CHEM-S2-Q03, source={TOPS - Test 03 (20 Nov 2022)}, paper=2, qno=3, type=numerical, subtopic={Electronic effects}, answer={8.00}, answer_src=key, difficulty=4, diff_parts={K1 S1 I1 C1 F1}, diff_note={concepts: count aromatic species among three drawings | steps: apply Huckel 4n + 2 for planar cyclic conjugated species; count in each structure}]
How many of these compounds / ions are aromatic?
\bfig{202.7}{TOPS03_p2_chem_s2_q03_1.png}
\bfig{212.8}{TOPS03_p2_chem_s2_q03_2.png}
\bfig{155.6}{TOPS03_p2_chem_s2_q03_3.png}
\end{bankq}
\begin{bankq}[id=TOPS03-P2-CHEM-S2-Q04, source={TOPS - Test 03 (20 Nov 2022)}, paper=2, qno=4, type=numerical, subtopic={Electronic effects}, answer={7.00}, answer_src=key, difficulty=5, diff_parts={K1 S2 I1 C1 F1}, diff_note={concepts: count species showing hyperconjugation | steps: alpha hydrogens adjacent to a sp2 carbon or carbocation or radical; check each structure; count}]
Find out the number of compounds / ions in which hyper conjugation is present?
\bfig{225.5}{TOPS03_p2_chem_s2_q04_1.png}

\bfig{231.7}{TOPS03_p2_chem_s2_q04_2.png}
\bfig{262.8}{TOPS03_p2_chem_s2_q04_3.png}
\end{bankq}
\begin{bankq}[id=TOPS03-P2-CHEM-S3-Q04, source={TOPS - Test 03 (20 Nov 2022)}, paper=2, qno=4, type=numerical, subtopic={Electronic effects}, answer={9}, answer_src=key, difficulty=4, diff_parts={K1 S1 I1 C1 F1}, diff_note={concepts: number of resonating structures of the given species | steps: draw all valid resonance forms by moving pi electrons and lone pairs; count}]
How many resonating structure are possible in
\bfig{81.6}{TOPS03_p2_chem_s3_q04_1.png}
\end{bankq}
\begin{bankq}[id=TOPS04-P1-CHEM-S1-Q03, source={TOPS - Test 04 (18 Dec 2022)}, paper=1, qno=3, type=multi, subtopic={Acidity and basicity}, answer={A,B,C,D}, answer_src=key, difficulty=4, diff_parts={K2 S1 I0 C1 F1}, diff_note={concepts: inductive effect, acidity | steps: effect of substituents; compare each pair}]
Correct acidity order
\begin{qoptions}
\opt{A}{$\mathrm{CF_3OH>CCl_3OH}$}
\opt{B}{$\mathrm{CHCl_3>CHF_3}$}
\opt{C}{\bfigi{30.1}{TOPS04_p1_chem_s1_q03_1.png} $>$ \bfigi{24.0}{TOPS04_p1_chem_s1_q03_C.png}}
\opt{D}{$\mathrm{Cl_3C{-}CO_2H<F_3C{-}CO_2H}$}
\end{qoptions}
\end{bankq}
\begin{bankq}[id=TOPS04-P1-CHEM-S1-Q04, source={TOPS - Test 04 (18 Dec 2022)}, paper=1, qno=4, type=multi, subtopic={Electronic effects}, answer={B,C,D}, answer_src=key, difficulty=5, diff_parts={K2 S1 I1 C1 F1}, diff_note={concepts: resonance, bond length | steps: bond order of C-O in each; compare with the reference}]
In which of the following compound bond length of C--O bond is greater then \bfigi{49.9}{TOPS04_p1_chem_s1_q04_1.png}
\begin{qoptions}
\opt{A}{\bfigi{47.8}{TOPS04_p1_chem_s1_q04_A.png}}
\opt{B}{\bfigi{47.8}{TOPS04_p1_chem_s1_q04_B.png}}
\opt{C}{\bfigi{47.8}{TOPS04_p1_chem_s1_q04_C.png}}
\opt{D}{\bfigi{50.4}{TOPS04_p1_chem_s1_q04_D.png}}
\end{qoptions}
\end{bankq}
\begin{bankpassage}[id=TOPS04-P1-CHEM-S2-G9-10, source={TOPS - Test 04 (18 Dec 2022)}, questions={9-10}]
Planar monocyclic ring containing $(4n+2)p$ electron (where $n=0,1,2,3\ldots$) are aromatic on applying this rule benzene and its derivative are aromatic.

Non benzenoid aromatic compound are those which contain ring that is not six membered such as \bfigi{24.0}{TOPS04_p1_chem_s2_para9_10_1.png}.

Cyclic planar conjugated compound which posses only $4np$ electron are antiaromatic in a reaction that reactant would tend to react with reagent that can result into aromatic product salt.
\end{bankpassage}
\begin{bankq}[id=TOPS04-P1-CHEM-S2-Q09, source={TOPS - Test 04 (18 Dec 2022)}, paper=1, qno=9, type=single, subtopic={Electronic effects}, answer={B}, answer_src=key, difficulty=4, diff_parts={K2 S1 I1 C1 F0}, diff_note={concepts: dipole moment, aromaticity | steps: dipole direction by aromatic stabilisation; check each}, group=TOPS04-P1-CHEM-S2-G9-10]
Which of the following is correct option regarding direction of dipole moment ?
\begin{qoptions}
\opt{A}{\bfigi{53.1}{TOPS04_p1_chem_s2_q09_A.png}}
\opt{B}{\bfigi{42.6}{TOPS04_p1_chem_s2_q09_B.png}}
\opt{C}{\bfigi{47.8}{TOPS04_p1_chem_s2_q09_C.png}}
\opt{D}{\bfigi{52.9}{TOPS04_p1_chem_s2_q09_D.png}}
\end{qoptions}
\end{bankq}
\begin{bankq}[id=TOPS04-P1-CHEM-S2-Q10, source={TOPS - Test 04 (18 Dec 2022)}, paper=1, qno=10, type=single, subtopic={Electronic effects}, answer={D}, answer_src=key, difficulty=4, diff_parts={K2 S1 I1 C1 F0}, diff_note={concepts: aromaticity, carbocation stability | steps: aromatic cation formation; reactivity with silver nitrate}, group=TOPS04-P1-CHEM-S2-G9-10]
Which of the following is correct ?
\begin{qoptions}
\opt{A}{\bfigi{47.8}{TOPS04_p1_chem_s2_q10_A.png} reacts immediately with $\mathrm{AgNO_3}$}
\opt{B}{\bfigi{47.2}{TOPS04_p1_chem_s2_q10_B.png} reacts immediately with $\mathrm{AgNO_3}$}
\opt{C}{\bfigi{156.0}{TOPS04_p1_chem_s2_q10_C.png}}
\opt{D}{Both (A) and (B)}
\end{qoptions}
\end{bankq}
\begin{bankq}[id=TOPS04-P2-CHEM-S1-Q07, source={TOPS - Test 04 (18 Dec 2022)}, paper=2, qno=7, type=single, subtopic={Electronic effects}, answer={C}, answer_src=key, difficulty=4, diff_parts={K2 S1 I1 C1 F0}, diff_note={concepts: heat of hydrogenation, stability | steps: stability of each alkene; compare}]
Which of the following is CORRECT :
\begin{qoptions}
\opt{A}{\bfigi{138.6}{TOPS04_p2_chem_s1_q07_A.png} hydrogenation}
\opt{B}{\bfigi{139.4}{TOPS04_p2_chem_s1_q07_B.png}}
\opt{C}{\bfigi{139.3}{TOPS04_p2_chem_s1_q07_C.png}}
\opt{D}{\bfigi{138.1}{TOPS04_p2_chem_s1_q07_D.png} Heat of hydrogenation}
\end{qoptions}
\end{bankq}
\begin{bankq}[id=TOPS04-P2-CHEM-S1-Q08, source={TOPS - Test 04 (18 Dec 2022)}, paper=2, qno=8, type=single, subtopic={Acidity and basicity}, answer={A}, answer_src=key, difficulty=4, diff_parts={K1 S2 I1 C1 F0}, diff_note={concepts: basicity | steps: hybridisation of nitrogen; lone pair availability; compare}]
Which compound is most basic -
\begin{qoptions}
\opt{A}{$\mathrm{CH_3{-}CH_2{-}NH_2}$}
\opt{B}{$\mathrm{CH_3{-}CH{=}NH}$}
\opt{C}{$\mathrm{CH_3{-}C{\equiv}N}$}
\opt{D}{\bfigi{24.0}{TOPS04_p2_chem_s1_q08_1.png}}
\end{qoptions}
\end{bankq}
\begin{bankq}[id=TOPS04-P2-CHEM-S1-Q09, source={TOPS - Test 04 (18 Dec 2022)}, paper=2, qno=9, type=single, subtopic={Acidity and basicity}, answer={C}, answer_src=key, difficulty=3, diff_parts={K1 S1 I1 C1 F0}, diff_note={concepts: acidity of enol hydrogens | steps: stability of each conjugate base; compare}]
Arrange marked hydrogen in decresing order of their acidic strength in vitamin C.
\bfig{69.0}{TOPS04_p2_chem_s1_q09_1.png}
\begin{qoptions}
\opt{A}{I $>$ II $>$ III}
\opt{B}{III $>$ II $>$ I}
\opt{C}{II $>$ III $>$ I}
\opt{D}{II $>$ I $>$ III}
\end{qoptions}
\end{bankq}
\begin{bankq}[id=TOPS04-P2-CHEM-S1-Q10, source={TOPS - Test 04 (18 Dec 2022)}, paper=2, qno=10, type=single, subtopic={Reactive intermediates}, answer={A}, answer_src=key, difficulty=4, diff_parts={K1 S2 I1 C1 F0}, diff_note={concepts: tautomerism | steps: alpha hydrogens; keto enol; check each}]
Which of the following compound will show tautomerism?
\begin{qoptions}
\opt{A}{\bfigi{32.0}{TOPS04_p2_chem_s1_q10_A.png}}
\opt{B}{$\mathrm{Me_3C{-}C{\equiv}N}$}
\opt{C}{\bfigi{55.3}{TOPS04_p2_chem_s1_q10_C.png}}
\opt{D}{\bfigi{44.8}{TOPS04_p2_chem_s1_q10_D.png}}
\end{qoptions}
\end{bankq}
\begin{bankq}[id=TOPS04-P2-CHEM-S2-Q05, source={TOPS - Test 04 (18 Dec 2022)}, paper=2, qno=5, type=numerical, subtopic={Acidity and basicity}, answer={2}, answer_src=solved, working={Release $\mathrm{CO_2}$ with $\mathrm{NaHCO_3}$: sulphonic acid, acetic acid, squaric acid, HCl, benzoic acid, cyclohexanecarboxylic acid, so $x=6$. Aromatic: benzenesulphonic acid, phenol, benzene, benzoic acid, so $y=4$. $a=2$}, difficulty=7, diff_parts={K2 S2 I2 C1 F1}, diff_note={concepts: acidity, aromaticity | steps: identify the acids; carbon dioxide release; aromatic compounds; difference}]
(a) \bfigi{28.2}{TOPS04_p2_chem_s2_q05_0.png}\hspace{0.5em}\bfigi{85.8}{TOPS04_p2_chem_s2_q05_1.png}
\bfig{139.9}{TOPS04_p2_chem_s2_q05_2.png}
\bfig{107.3}{TOPS04_p2_chem_s2_q05_3.png}
For the above compound identify the number

$a=x-y$

$x=$ Number of compounds which release $\mathrm{CO_2}$ on reaction with $\mathrm{NaHCO_3}$

$y=$ Number of compounds which are aromatic
\end{bankq}
\begin{bankq}[id=TOPS05-P1-CHEM-S3-Q04, source={TOPS - Test 05 (14 Jan 2023)}, paper=1, qno=4, type=numerical, subtopic={Acidity and basicity}, answer={4.00}, answer_src=key, difficulty=5, diff_parts={K1 S2 I1 C1 F1}, diff_note={concepts: compounds more acidic than a reference acid | steps: compare acid strengths of listed compounds by stability of the conjugate base; picric acid is stronger than phenol; sulphonic acid is stronger; count those stronger than the reference}]
Identify total number of compounds which are more acidic than \bfigi{60.7}{TOPS05_p1_chem_s3_q04_1.png}

\bfigi{65.8}{TOPS05_p1_chem_s3_q04_2.png} (2) Me--COOH \quad (3) Me--SO$_3$H \quad (4) Picric acid

\bfig{187.5}{TOPS05_p1_chem_s3_q04_3.png}
\end{bankq}
\begin{bankq}[id=TOPS05-P2-CHEM-S1-Q02, source={TOPS - Test 05 (14 Jan 2023)}, paper=2, qno=2, type=single, subtopic={Electronic effects}, answer={D}, answer_src=key, difficulty=2, diff_parts={K1 S1 I1 C0 F0}, diff_note={concepts: compounds that do not show tautomerism | steps: keto-enol tautomerism needs alpha hydrogens; check each structure}]
Which will not show tautomerism?
\begin{qoptions}
\opt{A}{\bfigi{64.2}{TOPS05_p2_chem_s1_q02_A.png}}
\opt{B}{\bfigi{76.8}{TOPS05_p2_chem_s1_q02_B.png}}
\opt{C}{\bfigi{60.5}{TOPS05_p2_chem_s1_q02_C.png}}
\opt{D}{\bfigi{65.8}{TOPS05_p2_chem_s1_q02_D.png}}
\end{qoptions}
\end{bankq}
\begin{bankq}[id=TOPS05-P2-CHEM-S2-Q07, source={TOPS - Test 05 (14 Jan 2023)}, paper=2, qno=7, type=multi, subtopic={Reactive intermediates}, answer={A}, answer_src=key, difficulty=3, diff_parts={K1 S1 I1 C0 F1}, diff_note={concepts: most stable carbocation | steps: compare stability by resonance and hyperconjugation and inductive effects}]
Which among the following carbocations is most stable?
\begin{qoptions}
\opt{A}{\bfigi{66.3}{TOPS05_p2_chem_s2_q07_A.png}}
\opt{B}{\bfigi{50.2}{TOPS05_p2_chem_s2_q07_B.png}}
\opt{C}{\bfigi{34.6}{TOPS05_p2_chem_s2_q07_C.png}}
\opt{D}{\bfigi{74.0}{TOPS05_p2_chem_s2_q07_D.png}}
\end{qoptions}
\end{bankq}
\begin{bankq}[id=TOPS05-P2-CHEM-S3-Q04, source={TOPS - Test 05 (14 Jan 2023)}, paper=2, qno=4, type=numerical, subtopic={Acidity and basicity}, answer={7}, answer_src=key, difficulty=5, diff_parts={K1 S2 I1 C1 F1}, diff_note={concepts: compounds soluble in aqueous NaOH | steps: acidic compounds react with NaOH: carboxylic acids phenols sulphonic acids; neutral compounds do not; count}]
How many compounds are soluble in aqueous NaOH?
\bfig{220.1}{TOPS05_p2_chem_s3_q04_1.png}
\bfig{217.2}{TOPS05_p2_chem_s3_q04_2.png}
(5) \bfigi{45.1}{TOPS05_p2_chem_s3_q04_3.png}
\end{bankq}

\section*{TMAS papers}
\begin{bankpassage}[id=TMAS02-P1-CHEM-S3-G3-4, source={TMAS - Test 02 (06 Aug 2023)}, questions={3-4}]
\par\noindent\hfill{\footnotesize$\text{Cyclopropyl benzene}\xrightarrow[\Delta]{\mathrm{H^+/KMnO_4}}\mathrm{(A)}\xrightarrow{\mathrm{PCl_5}}\mathrm{(B)}\xrightarrow[\mathrm{NH_4HCl}]{\mathrm{EtMgBr\ (1\ eq.)}}\mathrm{(C)}\xrightarrow{\mathrm{SeO_2}}\mathrm{(D)}\xrightarrow{\mathrm{HIO_4}}P+Q$}\hfill\null\par
In product $Q$, $4\pi$ electrons are involved in resonance
\end{bankpassage}
\begin{bankq}[id=TMAS02-P1-CHEM-S3-Q03, source={TMAS - Test 02 (06 Aug 2023)}, paper=1, qno=3, type=numerical, subtopic={Acidity and basicity}, answer={5.00}, answer_src=key, difficulty=6, diff_parts={K1 S2 I2 C1 F1}, diff_note={concepts: reaction sequence and unsaturation | steps: each step; product P; unsaturation}, group=TMAS02-P1-CHEM-S3-G3-4]
D.O.U. of $P$ is?
\end{bankq}
\begin{bankq}[id=TMAS02-P1-CHEM-S3-Q04, source={TMAS - Test 02 (06 Aug 2023)}, paper=1, qno=4, type=numerical, subtopic={Acidity and basicity}, answer={7.00}, answer_src=key, difficulty=5, diff_parts={K1 S2 I1 C1 F1}, diff_note={concepts: acidity comparison | steps: acidity of Q; compare each species; count}, group=TMAS02-P1-CHEM-S3-G3-4]
How many of the following are more acidic than $Q$
\bfig{334.6}{TMAS02_p1_chem_s3_q04_1.png}

$\mathrm{CH_3CH_2OH}$ , $\mathrm{HCl}$, $\mathrm{CF_3COOH}$
\end{bankq}

\section*{TNAS papers}
\begin{bankq}[id=TNAS01-P1-CHEM-S1-Q14, source={TNAS - Test 01 (05 Feb 2023)}, paper=1, qno=14, type=single, subtopic={Acidity and basicity}, answer={A}, answer_src=key, difficulty=2, diff_parts={K1 S1 I1 C0 F0}, diff_note={concepts: acidity order | steps: compare the groups in each option by inductive and resonance effects; find the correct one}]
Identify correct acidic strength order?
\begin{qoptions}
\opt{A}{\bfigi{136.4}{TNAS01_p1_chem_s1_q14_A.png}}
\opt{B}{\bfigi{35.4}{TNAS01_p1_chem_s1_q14_2.png}\ \bfigi{21.3}{TNAS01_p1_chem_s1_q14_3.png}\ $>$\ \bfigi{29.2}{TNAS01_p1_chem_s1_q14_4.png}}
\opt{C}{Me--C$\equiv$C--H $>$ MeOH $>$ $\mathrm{NH_3}$}
\opt{D}{\bfigi{137.4}{TNAS01_p1_chem_s1_q14_D.png}}
\end{qoptions}
\end{bankq}
\begin{bankq}[id=TNAS01-P1-CHEM-S1-Q15, source={TNAS - Test 01 (05 Feb 2023)}, paper=1, qno=15, type=single, subtopic={Reactive intermediates}, answer={C}, answer_src=key, difficulty=2, diff_parts={K1 S1 I1 C0 F0}, diff_note={concepts: carbocation stability | steps: compare hyperconjugation and resonance for three cations; order | I: resonance by adjacent pi bond}]
Choose the order that has the following carbocations correctly arranged with respect to increasing stability?
\bfig{69.3}{TNAS01_p1_chem_s1_q15_1.png}
\bfig{85.2}{TNAS01_p1_chem_s1_q15_2.png}
\bfig{66.6}{TNAS01_p1_chem_s1_q15_3.png}
\begin{qoptions}
\opt{A}{$1>2>3$}
\opt{B}{$1>3>2$}
\opt{C}{$2>1>3$}
\opt{D}{$3>2>1$}
\end{qoptions}
\end{bankq}
\begin{bankq}[id=TNAS01-P1-CHEM-S2-Q09, source={TNAS - Test 01 (05 Feb 2023)}, paper=1, qno=9, type=numerical, subtopic={Reactive intermediates}, answer={6}, answer_src=solved, working={Cations that rearrange: neopentyl type, isobutyl, the hydroxy cation, the tetralin cation, the secondary cation next to a tertiary C-H and the cyclopropyl substituted ring cation; the tertiary 3-methyl-3-pentyl cation does not}, difficulty=2, diff_parts={K1 S1 I0 C0 F1}, diff_note={concepts: carbocation rearrangement | steps: check adjacent hydrogen or alkyl shifts for each cation; count}]
Identify total number of carbocations which undergoes rearrangement reactions
\bfig{50.6}{TNAS01_p1_chem_s2_q09_1.png}
\bfig{50.6}{TNAS01_p1_chem_s2_q09_2.png}
\bfig{66.3}{TNAS01_p1_chem_s2_q09_3.png}
\bfig{43.1}{TNAS01_p1_chem_s2_q09_4.png}
\bfig{50.6}{TNAS01_p1_chem_s2_q09_5.png}
\bfig{60.8}{TNAS01_p1_chem_s2_q09_6.png}
\end{bankq}
\begin{bankq}[id=TNAS03-P1-CHEM-S1-Q13, source={TNAS - Test 03 (28 May 2023)}, paper=1, qno=13, type=single, subtopic={Acidity and basicity}, answer={B}, answer_src=key, difficulty=3, diff_parts={K1 S1 I1 C1 F0}, diff_note={concepts: stability order | steps: effects in each structure; compare}]
Select the order which is not correct among following :
\begin{qoptions}
\opt{A}{\bfigi{144.2}{TNAS03_p1_chem_s1_q13_A.png}}
\opt{B}{\bfigi{108.5}{TNAS03_p1_chem_s1_q13_B.png}}
\opt{C}{\bfigi{141.1}{TNAS03_p1_chem_s1_q13_C.png}}
\opt{D}{\bfigi{108.5}{TNAS03_p1_chem_s1_q13_D.png}}
\end{qoptions}
\end{bankq}
\begin{bankq}[id=TNAS03-P1-CHEM-S2-Q03, source={TNAS - Test 03 (28 May 2023)}, paper=1, qno=3, type=numerical, subtopic={Acidity and basicity}, answer={9}, answer_src=key, difficulty=5, diff_parts={K1 S2 I1 C1 F1}, diff_note={concepts: solubility in base | steps: acid strengths; reaction with NaOH and bicarbonate; sum}]
Compounds:
\bfig{158.5}{TNAS03_p1_chem_s2_q03_1.png}
(ix) $\mathrm{CH_3{-}CH_2{-}OH}$\\
$x$ = number of compounds soluble in NaOH (aq.)\\
$y$ = number of compounds soluble in $\mathrm{NaHCO_3}$ (aq.)\\
then $x+y=$ ?
\end{bankq}
\begin{bankq}[id=TNAS03-P2-CHEM-S3-Q05, source={TNAS - Test 03 (28 May 2023)}, paper=2, qno=5, type=numerical, subtopic={Acidity and basicity}, answer={10}, answer_src=key, difficulty=6, diff_parts={K1 S2 I2 C1 F1}, diff_note={concepts: aromaticity and basicity | steps: aromatic count; most basic; most acidic; sum}]
Consider the following compounds :
\bfig{208.6}{TNAS03_p2_chem_s3_q05_1.png}
Total number of aromatic compound is (R)\\
Serial number of most basic compound is (P)\\
Serial number of most acidic compound is (Q)\\
\textbf{Number of compound represented in bracket} then find value of $P+Q+R$?
\end{bankq}
\begin{bankq}[id=TNAS07-P2-CHEM-S1-Q06, source={TNAS - Test 07 (20 Aug 2023)}, paper=2, qno=6, type=multi, subtopic={Electronic effects}, answer={A,B,C,D}, answer_src=key, difficulty=4, diff_parts={K1 S2 I1 C0 F1}, diff_note={concepts: aromaticity of rings | steps: count pi electrons; planarity; conjugation}]
Among P, Q, R and S, the aromatic compound(s) is / are :
\bfig{168.0}{TNAS07_p2_chem_s1_q06_1.png}
\bfig{159.0}{TNAS07_p2_chem_s1_q06_2.png}
\begin{qoptions}
\opt{A}{P}
\opt{B}{Q}
\opt{C}{R}
\opt{D}{S}
\end{qoptions}
\end{bankq}

\section*{Internal Test papers}
\begin{bankq}[id=INT02-P1-CHEM-S1-Q01, source={Internal Test - Test 02 (07 Jul 2021)}, paper=1, qno=1, type=single, subtopic={Acidity and basicity}, answer={A}, answer_src=key, difficulty=2, diff_parts={K1 S2 I0 C0 F0}, diff_note={concepts: acid strength ordering | steps: compare the stability of the conjugate bases; inductive effect; hydrogen bonding}]
Arrange the following acids in decreasing order of acidity :\par
\bfig{201.4}{INT02_p1_chem_s1_q01_1.png}
\begin{qoptions}
\opt{A}{$X>Z>Y$}
\opt{B}{$X>Y>Z$}
\opt{C}{$Z>X>Y$}
\opt{D}{$Z>Y>X$}
\end{qoptions}
\end{bankq}
\begin{bankq}[id=INT02-P1-CHEM-S1-Q02, source={Internal Test - Test 02 (07 Jul 2021)}, paper=1, qno=2, type=single, subtopic={Electronic effects}, answer={D}, answer_src=key, difficulty=3, diff_parts={K1 S2 I1 C0 F0}, diff_note={concepts: stability of organic compounds | steps: resonance in each; hyperconjugation; strain; pick the most stable}]
Which of the following compounds is most stable ?
\begin{qoptions}
\opt{A}{\bfigi{54.7}{INT02_p1_chem_s1_q02_A.png}}
\opt{B}{\bfigi{54.7}{INT02_p1_chem_s1_q02_B.png}}
\opt{C}{\bfigi{68.9}{INT02_p1_chem_s1_q02_C.png}}
\opt{D}{\bfigi{75.1}{INT02_p1_chem_s1_q02_D.png}}
\end{qoptions}
\end{bankq}
\begin{bankq}[id=INT02-P1-CHEM-S3-Q13, source={Internal Test - Test 02 (07 Jul 2021)}, paper=1, qno=13, type=numerical, subtopic={Electronic effects}, answer={5}, answer_src=key, difficulty=3, diff_parts={K1 S2 I0 C0 F1}, diff_note={concepts: aromaticity | steps: planarity and cyclic conjugation of each species; 4n+2 electron rule; count}]
Find out number of aromatic compounds or ions from the following\par
\bfig{199.8}{INT02_p1_chem_s3_q13_1.png}
\end{bankq}
\begin{bankq}[id=INT02-P1-CHEM-S3-Q18, source={Internal Test - Test 02 (07 Jul 2021)}, paper=1, qno=18, type=numerical, subtopic={Electronic effects}, answer={5}, answer_src=key, difficulty=3, diff_parts={K1 S2 I0 C0 F1}, diff_note={concepts: benzylic hydrogens | steps: identify benzylic carbons in each structure; count hydrogens on them; total}]
Find out number of benzylic hydrogen in\par
\bfig{58.5}{INT02_p1_chem_s3_q18_1.png}
\end{bankq}
\begin{bankq}[id=INT03-P1-CHEM-S1-Q04, source={Internal Test - Test 03 (08 Sep 2021)}, paper=1, qno=4, type=single, subtopic={Acidity and basicity}, answer={B}, answer_src=key, difficulty=3, diff_parts={K1 S2 I1 C0 F0}, diff_note={concepts: strongest acid among a set | steps: substituents on the phenol; carboxylic acids; rank}]
Which of the following compounds is strongest acid?
\begin{qoptions}
\opt{A}{Carbolic acid}
\opt{B}{Picric acid}
\opt{C}{Salycilic acid}
\opt{D}{Formic acid}
\end{qoptions}
\end{bankq}

\section*{Nurture (General) papers}
\begin{bankq}[id=NUR04-P1-CHEM-S2-Q05, source={Nurture (General) - Test 04 (11 Feb 2024)}, paper=1, qno=5, type=multi, subtopic={Reactive intermediates}, answer={A,B,C}, answer_src=solved, working={Oxygen-stabilised cation, allylic against vinylic cation and nitro stabilisation of the carbanion hold; D is reversed}, difficulty=7, diff_parts={K2 S2 I2 C1 F1}, diff_note={concepts: stability of intermediates, reactive intermediates | steps: compare each pair; electronic effects; check options}]
Which of the following correctly represents the stability of reactive intermediate ?
\begin{qoptions}
\opt{A}{$\mathrm{CH_3-\overset{+}{C}H_2<CH_3O\overset{\oplus}{C}H_2}$}
\opt{B}{\bfigi{126.4}{NUR04_p1_chem_s2_q05_B.png}}
\opt{C}{\bfigi{90.5}{NUR04_p1_chem_s2_q05_C.png}}
\opt{D}{$\mathrm{(CH_3)_3\overset{\oplus}{C}<CH_3\overset{\oplus}{C}H_2}$}
\end{qoptions}
\end{bankq}
\begin{bankq}[id=NUR04-P1-CHEM-S2-Q08, source={Nurture (General) - Test 04 (11 Feb 2024)}, paper=1, qno=8, type=multi, subtopic={Electronic effects}, answer={A,B}, answer_src=solved, working={Cyclopropenone and tropone gain aromatic zwitterion character that evens out the C-O bond}, difficulty=7, diff_parts={K2 S2 I2 C1 F1}, diff_note={concepts: bond lengths in carbonyl compounds, electronic effects | steps: resonance contributors; aromatic character; compare}]
Which of the following will have C$=$O bond almost similar to C$-$O bond length ?
\begin{qoptions}
\opt{A}{\bfigi{36.5}{NUR04_p1_chem_s2_q08_A.png}}
\opt{B}{\bfigi{44.3}{NUR04_p1_chem_s2_q08_B.png}}
\opt{C}{\bfigi{39.5}{NUR04_p1_chem_s2_q08_C.png}}
\opt{D}{O$=$C$=$O}
\end{qoptions}
\end{bankq}
\begin{bankq}[id=NUR04-P1-CHEM-S3-Q15, source={Nurture (General) - Test 04 (11 Feb 2024)}, paper=1, qno=15, type=numerical, subtopic={Electronic effects}, answer={9}, answer_src=solved, working={The 2-naphthoxide ion has 3 Kekule structures with the charge on oxygen and 6 with the charge on C1, C3, C4a, C6, C8}, difficulty=7, diff_parts={K2 S2 I2 C1 F1}, diff_note={concepts: resonance structures of the naphthoxide ion, electronic effects | steps: charge on oxygen; positions of the carbanion; count the Kekule forms}]
The number of resonance structures for N is\par
\bfig{124.2}{NUR04_p1_chem_s3_q15_1.png}
\end{bankq}
\end{document}
